\documentclass[11pt]{article}

\usepackage[T1]{fontenc}
\usepackage[utf8]{inputenc}
\usepackage[margin=1in]{geometry}
\usepackage{lmodern}
\usepackage{microtype}
\usepackage{amsmath,amssymb,mathtools,amsthm}
\usepackage{aliascnt}
\usepackage{booktabs,tabularx,array,longtable,multirow}
\usepackage[section]{placeins}
\usepackage{float}
\usepackage{graphicx}
\usepackage{xcolor}
\usepackage{enumitem}
\usepackage{needspace}
\usepackage{url}
\usepackage[numbers,sort&compress]{natbib}
\usepackage[hidelinks]{hyperref}
\usepackage[nameinlink,noabbrev]{cleveref}

\allowdisplaybreaks
\setlength{\emergencystretch}{2em}
\raggedbottom
\setlist{leftmargin=*,topsep=3pt,itemsep=1.5pt,parsep=0pt,partopsep=0pt}

\definecolor{deepblue}{RGB}{23,69,124}
\hypersetup{
  colorlinks=true,
  linkcolor=deepblue,
  citecolor=deepblue,
  urlcolor=deepblue,
  pdftitle={A Symmetry-Quotient Reduction for Odd-Characteristic SNOVA: Exact Structure and Conditional Recovery Ledgers},
  pdfauthor={Justin Thaler, Kai Zhang, Scott Kominers, Quang Dao},
  pdfsubject={Cryptanalysis of odd-characteristic SNOVA and its Version 2.3 parameters},
  pdfkeywords={SNOVA, multivariate signatures, symmetric public matrices, forgery, NIST PQC}
}

\newtheorem{theorem}{Theorem}[section]
\newaliascnt{lemma}{theorem}
\newtheorem{lemma}[lemma]{Lemma}
\aliascntresetthe{lemma}
\newaliascnt{proposition}{theorem}
\newtheorem{proposition}[proposition]{Proposition}
\aliascntresetthe{proposition}
\newaliascnt{corollary}{theorem}
\newtheorem{corollary}[corollary]{Corollary}
\aliascntresetthe{corollary}
\newaliascnt{definition}{theorem}
\newtheorem{definition}[definition]{Definition}
\aliascntresetthe{definition}
\newaliascnt{remark}{theorem}
\newtheorem{remark}[remark]{Remark}
\aliascntresetthe{remark}

\crefname{theorem}{theorem}{theorems}
\Crefname{theorem}{Theorem}{Theorems}
\crefname{lemma}{lemma}{lemmas}
\Crefname{lemma}{Lemma}{Lemmas}
\crefname{proposition}{proposition}{propositions}
\Crefname{proposition}{Proposition}{Propositions}
\crefname{corollary}{corollary}{corollaries}
\Crefname{corollary}{Corollary}{Corollaries}
\crefname{definition}{definition}{definitions}
\Crefname{definition}{Definition}{Definitions}
\crefname{remark}{remark}{remarks}
\Crefname{remark}{Remark}{Remarks}

\newcommand{\F}{\mathbb F}
\newcommand{\Mat}{\operatorname{Mat}}
\newcommand{\Sym}{\operatorname{Sym}}
\newcommand{\rank}{\operatorname{rank}}
\newcommand{\Span}{\operatorname{span}}
\newcommand{\im}{\operatorname{im}}
\newcommand{\transpose}{\mathsf T}
\newcommand{\SNOVA}{\textsf{SNOVA}}
\newcommand{\AXN}{\mathsf{AXN}}
\newcommand{\attacktitle}{A Symmetry-Quotient Reduction for
  Odd-Characteristic \texorpdfstring{\SNOVA}{SNOVA}: Exact Structure and
  Conditional Recovery Ledgers}

\title{\attacktitle}
\author{%
  Justin Thaler\thanks{a16z crypto research and Georgetown University} \and
  Kai Zhang\thanks{MIT} \and
  Scott Kominers\thanks{Harvard University and a16z crypto research} \and
  Quang Dao\thanks{Carnegie Mellon University and LayerZero Labs}%
}
\date{31 July 2026}

\begin{document}
\maketitle

\begin{abstract}
\SNOVA{} is a multivariate signature scheme that compresses the public key of
Unbalanced Oil and Vinegar by representing variables and coefficients with
small matrices.  After wedge attacks weakened its characteristic-two design,
\SNOVA{} adopted odd-prime fields and symmetric public matrices.  This paper
shows that the symmetry creates a different weakness in the public verifier.

We analyze the common-column relation already used in
affine-column forgery attacks.  For a symmetric odd-characteristic \SNOVA{}
instance whose public matrix algebra has dimension $d$, symmetry reduces the
number of distinct restricted quadratic forms from $d^2$ to $d(d+1)/2$.
Thus sixteen forms become ten when $\ell=4$, and four become three when
$\ell=2$.  An invertible public change of output coordinates then separates
these quadratics from the remaining affine constraints and target checks.
For every published Version~2.3 parameter shape that passes the stated public
rank checks, this gives an exact quadratic residual system on a public affine
space.  The reduction permits affine slices of any dimension and preserves
every verifier equation, rejection rule, and final verification check.

The recovery estimates use Polynomial XL.  For the three $\ell=2$ shapes, we
prove full affine coefficient support on transverse slices.  For $\ell=4$,
we prove an exact graph-complement realization of all ten homogeneous
channels on a zero-offset slice.  Positive-rank offsets require a separate
publicly checkable compatibility premise whose frequency under official key
generation is unproved.  Under the ordinary PXL assumptions and an
instancewise extraction premise, the support-constrained route lies below
the Level III and V reference exponents but not Level I.  A second,
manuscript-specific structured PXL premise gives conditional estimates below
the reference exponents for all nine shapes.  The lowest unit-leading gate
exponents for $\omega=2.81$ are $128.833$, $187.030$, and $245.042$ for
$\ell=2$, and $134.038$, $182.006$, and $230.444$ for $\ell=4$.  These
figures omit auxiliary work and are not wall-clock or peak-memory bounds.  We
do not claim a production-size solve or forgery.
\end{abstract}

\section{Introduction}\label{sec:intro}

NIST's Additional Digital Signatures process seeks signature schemes with
security and performance profiles that complement the standards selected in
its first post-quantum process.  \SNOVA{} is one of the nine candidates
advanced to the third round in May 2026~\cite{NISTRound3,NISTIR8610}.  It
belongs to the Unbalanced Oil and Vinegar line of multivariate signatures,
where verification evaluates a public quadratic map and signing uses a
hidden central map with an oil-and-vinegar structure
~\cite{KipnisPatarinGoubin1999,SNOVA23}.

\SNOVA{} obtains compact public keys by representing variables and
coefficients with small matrices and generating many public forms from a
smaller structured family.  That structure has led to key-recovery and direct
forgery attacks.  The former expose the hidden UOV structure, while Beullens
uses a public relation among signature columns.  Later work improves both
lines through group actions, extension-field lifting, and
multihomogeneous solving
~\cite{IkematsuAkiyama2024,LiDing2024,Beullens2024SNOVA,Cabarcas2025SNOVA,FurueEtAl2025,NakamuraEtAl2025}.
These works analyze the original or Round~2 construction over $\F_{16}$.

The next revision answered Ran's wedge attack, which uses the alternating
polar forms available in characteristic two.  Bros et al. adapt the method to
SNOVA's block structure~\cite{Ran2026,BrosEtAl2026SNOVAWedge}.  The \SNOVA{}
team responded with odd-prime fields and symmetric public quadratic matrices,
while warning that symmetrization could create new attacks
~\cite{WangEtAl2025Round2Security}.  NIST later described this change as a
response to the wedge attack and did not regard \SNOVA{} as stable
~\cite{NISTIR8610}.

This change entered the specification in Version~2.1, not Version~2.3.
Version~2.3 subsequently added rectangular signatures and revised parameter
sets~\cite{SNOVA23}.  Our quotient theorem concerns the symmetric
odd-characteristic construction and is not specific to those nine parameter
sets.  The exact rank and dimension analysis covers all nine published
$q=19$ Version~2.3 shapes.  The audited numerical PXL ledger covers both
extension degrees, with separate premises for the support-constrained and
structured $\ell=4$ routes.

Prior attacks already cover odd characteristic.  Jin et al. give a UOV
key-recovery attack in arbitrary characteristic~\cite{JinEtAl2026}.  Ran gives
an odd-characteristic construction using symmetric forms, but SNOVA's added
structure requires separate analysis~\cite{Ran2026}.  Ding, Li, and Tao give a
wedge framework valid in any characteristic, while their concrete SNOVA
estimates cover the even-characteristic Round~2 parameters
~\cite{DingLiTao2026}.  None combines the symmetric public matrices in the
current $q=19$ construction with the affine-column forgery relation.

This brings us to the question addressed in this paper.  What happens when an
affine-column attacker imposes the same common-column relation on the new
symmetric verifier?  We prove that symmetry makes pairs of restricted
quadratic forms identical.  An invertible public change of output coordinates
then separates these quadratics from the remaining affine constraints and
target checks.  For every published $q=19$ parameter shape that passes the
stated public rank checks, this leaves a quadratic system on a public affine
space.  The attacker may choose a square, overdetermined, or underdetermined
affine slice without dropping a verifier equation.

This statement concerns the restricted verifier, not the unrestricted
matrix-valued map.  Let $A=\F_q[S]$ have dimension $d$, and let
$\Sigma_\xi=I_n\otimes\xi$ for $\xi\in A$.  After all signature columns are
expressed through one stacked column $u$, symmetry gives
$q_{i,\xi,\eta}(u)=q_{i,\eta,\xi}(u)$.  The $d^2$ ordered labels therefore
factor through $\Sym^2(A)$, with the exact dimension change
\begin{equation}
 d^2\ \longrightarrow\ \binom{d+1}{2},
 \qquad 16\ \longrightarrow\ 10\ \text{ for }d=4,
 \qquad 4\ \longrightarrow\ 3\ \text{ for }d=2.
 \label{eq:intro-duplication}
\end{equation}
The full verifier relates the two ordered features by matrix transposition;
only the common-column restriction makes them the same scalar polynomial.
The identification summarized in \Cref{eq:intro-duplication} is exact and
basis-independent, but deliberately local to the attacker-imposed
restriction (\Cref{thm:symmetry-quotient}).

\subsection{Our results}

\paragraph{An exact symmetry quotient and public reduction.}
After the common-column restriction, every symmetric public base form factors
through $\Sym^2(A)$.  This loss occurs before the outer public mixing, so later
linear combinations cannot restore the missing forms.  Public row reduction
then separates the surviving quadratics from the affine conditions and target
checks.  For every published parameter shape that passes the stated public
rank checks, this gives an exact quadratic system on a public affine space.
Every affine restriction preserves all residual equations, and every
recovered root is reconstructed and checked by the original verifier
(\Cref{thm:symmetry-quotient,prop:exact-affine,thm:exact-all-nine}).

\paragraph{Conditional PXL estimates for $\ell=2$.}
For $\ell=2$, we prove full affine coefficient support on every transverse
slice.  A uniform target and affine-coset trial then has an exact first moment,
and the spectrum theorem gives a lower bound on the probability of an
accepted root with full column rank.  We apply the five components from
Section~4.2 of PXL under its ordinary MQ assumptions and a separate
instancewise extraction premise.  For $\omega=2.81$, the resulting
unit-leading gate exponents are $128.833$, $187.030$, and $245.042$ across
Levels I, III, and V (\Cref{thm:pxl-l2-estimates}).  The NIST reference
exponents are $143$, $207$, and $272$~\cite{NISTSecurityCriteria2026}.
Auxiliary operations remain explicit and unpriced, so these figures are not
complete operation bounds.

\paragraph{The stronger $\ell=4$ quotient and two PXL routes.}
For $\ell=4$, symmetry removes six of every sixteen ordered quadratic forms.
After splitting the four Frobenius blocks, we prove that the fixed coefficient
map to the four diagonal and six cross channels is onto exactly when all four
maps to the vinegar quotients are injective
(\Cref{thm:pxl-fixed-map-surjectivity}).  The official graph complement gives
an exact zero-offset realization of this homogeneous support criterion.  For
positive-rank offsets, we state exact public compatibility checks as
$\mathsf H_{\rm off}$; their frequency under the official distribution is
unproved.  Under $\mathsf H_{\rm off}$ and the instancewise PXL route premise,
the support-constrained estimates cross the Level III and V reference
exponents but miss Level I.  A separate premise $\mathsf H_{\rm struct}$
asserts the required Hilbert series, Macaulay ranks, degree, and extraction
behavior directly for the structured channel family.  Under that premise,
the PXL estimates cross the reference exponents for all six $\ell=4$ shapes
(\Cref{thm:pxl-fixed-map-surjectivity,thm:pxl-all-nine-structured}).

\paragraph{Design implications.}
We derive necessary output-dimension conditions for redesigns that retain the
common-column architecture.  These conditions are not sufficient for security
(\Cref{thm:upgrade-combined-floor,thm:upgrade-zero-offset-floor}).

\subsection{How to read the result}

The paper separates the public reduction, the existence of acceptable roots,
and the behavior of the algebraic solver.  These statements have different
logical status.  The following table records the distinction used throughout
the paper.

\begin{table}[H]
\centering
\caption{Logical status of the attack.  An exact statement is a deterministic
claim about the supplied public key after its stated public checks pass.}
\label{tab:logical-status}
\small
\renewcommand{\arraystretch}{1.11}
\begin{tabularx}{0.98\textwidth}{@{}>{\raggedright\arraybackslash}p{0.25\textwidth}>{\raggedright\arraybackslash}p{0.19\textwidth}X@{}}
\toprule
Part & Status & Meaning\\
\midrule
Column restriction and symmetry quotient & Exact & Public substitution followed by an algebraic identity.\\
Output split and affine elimination & Exact after public rank checks & Invertible row reduction preserves every verifier equation.\\
Affine coefficient support for $\ell=2$ & Exact on transverse slices & The coefficient map is onto.  This does not imply uniform coefficients or semi-regularity.\\
Frequency of spectrum events & Distributional, or certified for a fixed key & The paper uses an idealized XOF-transcript model; a fixed key instead requires a spectrum certificate.\\
PXL regularity model & Conditional & The solving degree and matrix ranks use the ordinary MQ assumptions from PXL Section~4.2.\\
Instancewise PXL route & Conditional & The actual restricted system must follow the predicted ranks, and the added extraction step must recover a root when one exists.\\
$\ell=4$ homogeneous fixed-map criterion & Exact after splitting & Surjectivity holds exactly under four quotient-injectivity conditions, and the official graph complement realizes the zero-offset case.\\
$\ell=4$ positive-rank offsets & Conditional, with exact public checks & $\mathsf H_{\rm off}$ requires $\rank(C|_{U_{\rm vin}})=\lambda$ and $A\ker(C|_{U_{\rm vin}})=U_{\rm vin}$; its official frequency is unproved.\\
Structured $\ell=4$ PXL route & Conditional & $\mathsf H_{\rm struct}$ directly assumes the required Hilbert series, degree, Macaulay ranks, and extraction behavior.  It does not assert ordinary coefficient support.\\
PXL arithmetic & Unit-leading estimate & The five published components are aggregated conservatively.  Auxiliary work and storage traffic remain unpriced.\\
Candidate acceptance & Exact & Every candidate is reconstructed and checked by the original verifier.\\
Practical time and memory & Not established & No production solve is given, and the analysis does not bound wall-clock time or peak memory.\\
\bottomrule
\end{tabularx}
\end{table}

\subsection{Attack overview}

In the hidden UOV coordinates, \SNOVA{} contains no product of two oil
blocks.  The signer first chooses the vinegar blocks, after which the signing
equations are linear in the oil entries.  A secret invertible change of
variables hides this division, so the public verifier presents a structured
quadratic map rather than the linear system used by the signer
(\Cref{sec:scheme}).

The attack seeks a direct public preimage, not the hidden oil space.  It starts
from the common-column relation used in earlier forgery attacks.  Inside every
matrix-valued signature block, the attacker writes all columns as public
affine functions of one stacked column.  This selects a public subset of
possible signatures; it is not an identity of the unrestricted verifier
(\Cref{sec:public-symmetry}).

After this restriction, every ordered pair of public matrix multipliers gives
a scalar quadratic polynomial.  Symmetry makes the polynomials for the two
orders equal.  One group of sixteen ordered pairs therefore contributes at
most ten distinct quadratics when $\ell=4$, while one group of four
contributes at most three when $\ell=2$.  The equality holds before the outer
public map combines these forms, so later linear combinations cannot restore
the missing quadratic directions (\Cref{thm:symmetry-quotient}).

The attacker next computes an invertible public change of verifier outputs.
Some new coordinates contain the surviving quadratic forms.  The remaining
coordinates are affine conditions on the signature and, when necessary,
consistency conditions on the target.  Solving the affine equations gives a
public affine space.  Restricting this space to as many variables as surviving
quadratics gives a square system, but this is only one possible restriction.
Any affine slice preserves every residual equation
(\Cref{prop:exact-affine,cor:exact-affine-restriction}).

For $\ell=2$, the attacker samples a transverse affine coset and keeps all
residual equations.  The exact support theorem shows that the native
coefficient map reaches every ordinary affine MQ coefficient tuple on this
slice.  The root theorem bounds the chance that one target and coset trial
contains an accepted root with full column rank.  The conditional solver then
applies PXL to the complete overdetermined system
(\Cref{sec:pxl-l2-slice,sec:pxl-l2-recovery}).

For $\ell=4$, the exact quotient leaves ten channels per base form.  The
split fixed-map theorem identifies the precise condition under which these
homogeneous channels have full support, and the official graph complement
realizes it exactly at zero offset.  The positive-rank route uses the public
compatibility premise $\mathsf H_{\rm off}$.  Its support-constrained PXL
rows cross Levels III and V but not Level I.  The cheaper structured route
states its Hilbert-series, degree, Macaulay-rank, and extraction premise
directly and gives conditional rows below the reference exponents for all
six $\ell=4$ shapes (\Cref{sec:pxl-l4-support,sec:pxl-l4-structured}).

Finally, every recovered root is converted back into a candidate signature,
checked against the public format rule, and passed to the original verifier.
A failed rank test, target check, solve, or final verification causes a
restart.  The final check makes the attack one-sided: it can fail to return a
signature, but it never returns an invalid one.

\section{The \SNOVA{} signature scheme}\label{sec:scheme}

This section defines the part of Version~2.3 used in the attack.  Fix
parameters $(v,o,q,\ell,r)$ and put $n=v+o$.  A vector of private or public
variables has the form
\[
 X=(X_1,\ldots,X_n)\in\Mat_{\ell\times r}(\F_q)^n.
\]
In the private coordinates, $X_1,\ldots,X_v$ are the vinegar blocks and
$X_{v+1},\ldots,X_n$ are the oil blocks.  The central map has $o$ outputs in
$\Mat_{r\times\ell}(\F_q)$, so both its oil input and its output contain
$or\ell$ field elements~\cite{KipnisPatarinGoubin1999,SNOVA23}.

\paragraph{The hidden central map.}
The UOV structure excludes every product of two oil variables.  Write
$[n]=\{1,\ldots,n\}$.  The allowed ordered pairs are
\begin{equation}
 \Omega=[n]^2\setminus\{v+1,\ldots,n\}^2.
 \label{eq:scheme-omega}
\end{equation}
Let $S\in\Mat_\ell(\F_q)$ be the public symmetric matrix with irreducible
characteristic polynomial, and let
\[
 A=\F_q[S]\cong\F_{q^\ell}.
\]
Every element of $A$ is a symmetric matrix, and all elements of $A$ commute.
Version~2.3 uses
\[
 m_1=\left\lceil\frac{or}{\ell}\right\rceil,
 \qquad N_\alpha=r\ell+r
\]
base forms and outer summands.  For each $h\in[m_1]$, the secret base form is
an $n$ by $n$ block matrix $[F_h]=(F_{h,jk})$, where
$F_{h,jk}\in\Mat_\ell(\F_q)$ and $F_{h,jk}=0$ whenever
$(j,k)\notin\Omega$.

We write $L_{i,\alpha}$ and $R_{i,\alpha}$ for the matrices called
$A_{i,\alpha}$ and $B_{i,\alpha}$ in the specification.  Their dimensions are
$r$ by $r$ and $r$ by $\ell$, respectively.  The other public factors
$Q_{i,\alpha,1},Q_{i,\alpha,2}$ are nonzero elements of $A$.  Put
$h(i,\alpha)=1+((i-1+\alpha)\bmod m_1)$.  The $i$th component of the hidden
central map is
\begin{equation}
 \widetilde F_i(X)=
 \sum_{\alpha=0}^{N_\alpha-1} L_{i,\alpha}
 \left(
   \sum_{(j,k)\in\Omega}
   X_j^\transpose
   Q_{i,\alpha,1}F_{h(i,\alpha),jk}Q_{i,\alpha,2}X_k
 \right)R_{i,\alpha},
 \qquad i\in[o].
 \label{eq:scheme-central}
\end{equation}
The inner sum is an $r$ by $r$ matrix, and the full output is an $r$ by
$\ell$ matrix.  Since \Cref{eq:scheme-central} has no product of two oil
blocks, fixing the vinegar blocks makes it linear in every entry of the oil
blocks.

\paragraph{The secret change of variables and the public map.}
The secret matrix $T^{12}=(T_{ab})$ has $v$ by $o$ entries in $A$.  It defines
the invertible change of variables $X=T(U)$ by
\begin{equation}
 \begin{aligned}
  X_a&=U_a-\sum_{b=1}^o T_{ab}U_{v+b}, &&1\le a\le v,\\
  X_{v+b}&=U_{v+b}, &&1\le b\le o.
 \end{aligned}
 \label{eq:scheme-secret-change}
\end{equation}
Thus $T$ hides the private vinegar and oil coordinates while preserving the
matrix shape of each block.  For every base form, key generation computes
\begin{equation}
 [P_h]=[T]^\transpose[F_h][T].
 \label{eq:scheme-base-public}
\end{equation}
The commutativity of $A$ then gives the public quadratic map
\begin{equation}
 \begin{split}
 \widetilde P_i(U)
 &=\sum_{\alpha=0}^{N_\alpha-1} L_{i,\alpha}
 \left(
   \sum_{j,k=1}^n
   U_j^\transpose
   Q_{i,\alpha,1}P_{h(i,\alpha),jk}Q_{i,\alpha,2}U_k
 \right)R_{i,\alpha}\\
 &=\widetilde F_i(T(U)).
 \end{split}
 \label{eq:scheme-public-map}
\end{equation}
For odd $q$, the public base forms are symmetric.  In block notation, this
means
\begin{equation}
 P_{h,kj}=P_{h,jk}^\transpose.
 \label{eq:scheme-public-symmetry}
\end{equation}
Equivalently, expanding all block entries over $\F_q$ gives a symmetric
$n\ell$ by $n\ell$ matrix $P_h$.  This is the symmetry used by our attack.

The block equations also explain the public key compression.  Partition the
base forms according to the vinegar and oil coordinates:
\[
 [F_h]=
 \begin{bmatrix}F_h^{11}&F_h^{12}\\F_h^{21}&0\end{bmatrix},
 \qquad
 [T]=
 \begin{bmatrix}I&-T^{12}\\0&I\end{bmatrix},
 \qquad
 [P_h]=
 \begin{bmatrix}P_h^{11}&P_h^{12}\\P_h^{21}&P_h^{22}\end{bmatrix}.
\]
The public seed determines the outer matrices, the $Q$ matrices, and
$P_h^{11},P_h^{12},P_h^{21}$.  The private seed determines $T^{12}$.  Key
generation reconstructs the hidden blocks and the remaining public block by
\begin{equation}
 \begin{aligned}
  F_h^{11}&=P_h^{11},
  &F_h^{12}&=P_h^{12}+P_h^{11}T^{12},\\
  F_h^{21}&=P_h^{21}+(T^{12})^\transpose P_h^{11},
  &P_h^{22}&=-(T^{12})^\transpose F_h^{12}-P_h^{21}T^{12}.
 \end{aligned}
 \label{eq:scheme-keygen-blocks}
\end{equation}
Thus the compressed public key is the public seed together with the blocks
$P_h^{22}$ for $h\in[m_1]$.  The secret key contains the seeds needed to
regenerate $T$ and the hidden forms~\cite{SNOVA23}.

\paragraph{Signing and verification.}
For a message $M$ and a 16-byte salt, the target is
\begin{equation}
 Y=H\bigl(s_{\rm public}\mathbin\|\operatorname{SHAKE256}(M,64)
          \mathbin\|\mathrm{salt}\bigr)
 \in\Mat_{r\times\ell}(\F_q)^o.
 \label{eq:scheme-target}
\end{equation}
The signer derives the vinegar blocks from the private seed, $M$, the salt,
and an attempt counter.  Write $V=(X_1,\ldots,X_v)$ for the chosen vinegar
blocks and let $z$ list the $or\ell$ entries of the oil blocks.  By
\Cref{eq:scheme-omega}, the signing equation is a square linear system
\begin{equation}
 G_Vz=\operatorname{vec}(Y)-c_V,
 \label{eq:scheme-signing-system}
\end{equation}
where $c_V=\operatorname{vec}(\widetilde F(V,0))$ and $G_V$ is the coefficient
matrix of the terms containing one vinegar block and one oil block.  If $G_V$
is singular, the signer chooses a new attempt.  Otherwise it solves for the
oil blocks, sets $U=T^{-1}(X)$, and retries if $U$ violates the public format
rule.  The signature is $U\mathbin\|\mathrm{salt}$.

The verifier expands the public key, checks the same format rule, recomputes
$Y$, and accepts exactly when
\begin{equation}
 \widetilde P(U)=Y.
 \label{eq:scheme-verification}
\end{equation}
Version~2.1 introduced the odd-prime fields and
\Cref{eq:scheme-public-symmetry}.  Version~2.3 retained that symmetry and
allowed $r\ne\ell$, which gives rectangular signature blocks.

\paragraph{Column form of the public verifier.}
For each column $c\in\{0,\ldots,r-1\}$, stack the $c$th columns of the public
signature blocks as
\begin{equation}
 u_c=\bigl(U_1[:,c]^\transpose,\ldots,U_n[:,c]^\transpose\bigr)^\transpose
 \in\F_q^{n\ell}.
 \label{eq:scheme-stacked-columns}
\end{equation}
Using the scalar expansion of $[P_h]$ and
$\Sigma_\xi=I_n\otimes\xi$, each scalar entry inside
\Cref{eq:scheme-public-map} is a linear combination of terms
\begin{equation}
 u_c^\transpose\Sigma_\xi P_h\Sigma_\eta u_d,
 \qquad \xi,\eta\in A.
 \label{eq:scheme-power-pair}
\end{equation}
The outer matrices $L_{i,\alpha}$ and $R_{i,\alpha}$ combine these terms by a
public linear map.  Thus SNOVA's matrix structure affects the scalar verifier
through the ordered left and right algebra multipliers in
\Cref{eq:scheme-power-pair}.  Section~4 imposes the attacker's column relation
and determines which of these ordered terms remain distinct.

\section{Prior attacks and mathematical background}\label{sec:background}

\subsection{SNOVA's attack and revision history}

The first line of \SNOVA{} cryptanalysis attacks the hidden UOV structure.
Ikematsu and Akiyama analyze key recovery on the scalar expansion of the core
map, and Li and Ding sharpen the interpretation as a UOV system with
$\ell v$ vinegar variables and $\ell o$ oil variables over the base field
~\cite{IkematsuAkiyama2024,LiDing2024}.  Nakamura, Tani, and Furue instead
lift the same structure to a smaller UOV instance over $\F_{q^\ell}$
~\cite{NakamuraEtAl2025}.  These attacks seek the hidden oil space or an
equivalent signing key.  They concern the original characteristic-two
construction and do not use the symmetric public matrices introduced later.

Beullens exposes a different layer.  He rewrites \SNOVA{} as a structured
MAYO-style whipping of a bilinear map, chooses an affine relation among the
signature columns, and searches for rank-deficient combinations in the outer
coefficient map~\cite{Beullens2024SNOVA}.  The result is a forgery attack,
including weak-key instances much cheaper than the generic case.  The Round~2
changes vary the outer coefficient matrices and mitigate the largest rank
drops, but they retain the affine-column attack surface.  Our attack starts
from this public restriction and acts one layer earlier: symmetry identifies
the restricted power-pair polynomials before the outer map sees them.

Cabarcas et al. and Furue et al. improve the algebraic solving of the systems
arising in the earlier attacks~\cite{Cabarcas2025SNOVA,FurueEtAl2025}.  The
former exploits stability under a commutative group action and adapts the
result to reconciliation and direct forgery.  The latter writes the
extension-field transformation explicitly and applies the resulting
multihomogeneous systems to intersection and rectangular-MinRank key recovery.
Their estimate for the old Round~2 $(q,v,o,\ell)=(16,29,6,5)$ set is below the
claimed Level-V threshold, subject to their stated solving-degree heuristic.
Both papers study the characteristic-two construction.  Neither derives a
loss of independent verifier equations from the later symmetric public
matrices.

Ran's wedge attack begins from the alternating polar forms of
characteristic-two UOV and recovers the hidden oil space through exterior
algebra~\cite{Ran2026}.  His latest version also gives a different construction
for odd characteristic, but states that SNOVA's added structure prevents the
analysis from carrying over directly.  Bros et al. specialize the
characteristic-two method to \SNOVA{} and give conjectural estimates below the
claimed security levels for every Round~2 parameter set
~\cite{BrosEtAl2026SNOVAWedge}.  The \SNOVA{} team's later analysis replaces
the conjectural kernel formula with a generating function and finds six
Round~2 sets below target~\cite{TsengEtAl2026Wedge}.

Ding, Li, and Tao place these attacks in a common Macaulay-matrix framework and
extend it to multihomogeneous systems in any characteristic
~\cite{DingLiTao2026}.  Their formulation applies to SNOVA's hidden oil-space
problem, but their concrete SNOVA analysis is expressly limited to the
even-characteristic Round~2 parameters.  It does not analyze the symmetric
public matrices or the nine $q=19$ Version~2.3 shapes.

Jin et al. give the main prior attack that genuinely crosses the
characteristic boundary~\cite{JinEtAl2026}.  Their $p$-reduced symmetric
algebra recovers Ran's exterior algebra at $p=2$ and defines an XL
key-recovery attack for arbitrary characteristic.  Their concrete
odd-characteristic application is to the twelve QR-UOV parameter sets, where
the attack does not beat reconciliation on the recommended parameters.  The
method targets the hidden UOV oil space.  It does not analyze the symmetric
public-matrix encoding of odd-characteristic \SNOVA{}, and it does not produce
the verifier-preserving direct-preimage reduction used here.

The distinction between these attacks is summarized in
\Cref{tab:prior-attacks}.  The exact new statement in this paper is the
$\Sym^2(A)$ factorization after the public common-column restriction and the
resulting invertible change of verifier outputs.  The direct-forgery cost
estimates built on that reduction remain conditional on their stated
root-count and solver hypotheses.

\begin{table}[H]
\centering
\caption{Scope of the closest attacks.  ``Output'' is what a successful attack
recovers; it is not a claim that every cited estimate is practical.}
\label{tab:prior-attacks}
\footnotesize
\renewcommand{\arraystretch}{1.08}
\begin{tabularx}{0.98\textwidth}{@{}>{\raggedright\arraybackslash}p{0.19\textwidth}>{\raggedright\arraybackslash}p{0.16\textwidth}>{\raggedright\arraybackslash}p{0.25\textwidth}>{\raggedright\arraybackslash}X@{}}
\toprule
Work & Targeted regime & Main mechanism & Output and relation to this paper\\
\midrule
Ikematsu--Akiyama; Li--Ding; Nakamura--Tani--Furue
& Original / early $q=16$
& Scalar or extension-field UOV interpretation
& Oil-space or equivalent-key recovery; no odd-$q$ public symmetrization.\\
Beullens; Cabarcas et al.; Furue et al.
& Original and Round~2 $q=16$
& Affine-column rank drop, group action, and multihomogeneous solving
& Direct forgery or key recovery; supplies the public restriction that this paper composes with symmetry.\\
Ran; Bros et al.; Tseng et al.; Ding--Li--Tao
& Arbitrary-characteristic UOV frameworks; concrete SNOVA estimates for Round~2 $q=16$
& Exterior or symmetric forms, Macaulay matrices, and the wedge-map kernel
& Oil-space recovery; no concrete analysis of the symmetric $q=19$ SNOVA verifier.\\
Jin et al.
& UOV families in arbitrary characteristic; concrete odd-$q$ QR-UOV instances
& $p$-reduced symmetric-algebra XL
& Oil-space recovery; no analysis of symmetric odd-$q$ SNOVA public matrices.\\
This paper
& All nine $q=19$ Version~2.3 shapes
& Common-column restriction, exact $\Sym^2(A)$ quotient, and public output split
& Conditional direct preimages checked by the original verifier; no hidden-key recovery claim.\\
\bottomrule
\end{tabularx}
\end{table}

The complementary generic-MQ and partitioning literature includes
Hashimoto's underdefined optimizer, Just Guess, generalized partitioning, and
recent F4 engineering
~\cite{Hashimoto2023,MayOstuzziRessler2026,AsanumaEtAl2026,SakataTakagi2026}.
Polynomial XL performs part of the Macaulay elimination over a polynomial
ring before evaluating selected variables~\cite{FurueKudo2024PXL}.  Its
Section~4.2 estimates assume affine semi-regularity, projected
semi-regularity, and predicted Macaulay ranks.  Our solver model applies those
standard premises to the exact quotient system and states the added
instancewise extraction premise separately.

\subsection{MQ systems, roots, and affine slices}
An MQ system over $\F_q$ is a map $g:\F_q^N\to\F_q^K$ whose coordinates
have total degree at most two, together with a target $z\in\F_q^K$.  Linear
and constant terms are allowed.  A root is an $x$ satisfying $g(x)=z$.
Underdetermination makes roots more plentiful but does not make the system
linear or automatically easy.

An affine space is a translate $a+V$ of a linear subspace.  Restricting the
search to an affine slice $a+L\subseteq a+V$ turns the problem into a system
in $h=\dim L$ coordinates.  Three events must be separated:
\begin{enumerate}[label=(\roman*)]
\item the target--slice pair contains a base-field root;
\item it contains an accepted root at which a chosen Jacobian has full rank;
\item the chosen algebraic solver recovers every such root within its stated
      work bound.
\end{enumerate}
The first two are properties of the map and slice.  The third is a solver
statement.  The later recovery analysis proves them separately.

\section{From the public verifier to an affine quadratic system}
\label{sec:public-symmetry}\label{sec:affine-reduction}

This section starts from the column form of the public verifier in
\Cref{eq:scheme-stacked-columns,eq:scheme-power-pair} and gives the exact attack
as one sequence of public operations.  The attacker imposes the common-column
relation.  Symmetry then identifies restricted quadratic forms.  An invertible
change of verifier outputs separates the remaining quadratic equations from
the affine ones.  Once the public rank checks pass, the attacker may restrict
the resulting affine space to any chosen dimension while keeping every
residual equation.  This section makes no claim that the resulting system has
a root or that a solver can find one.

\subsection{Parameters and public forms}
Retain the notation of Section~2 and put
\begin{equation}
 m_1=\left\lceil\frac{or}{\ell}\right\rceil,
 \qquad M=or\ell,
 \qquad N_0=n\ell,
 \qquad K=m_1\binom{\ell+1}{2}.                    \label{eq:parameters}
\end{equation}
The ceiling follows the pinned Version~2.3 reference
implementation~\cite{SNOVA23}.  The theorem below proves that $K$ is the
maximum number of distinct restricted quadratic forms left by symmetry.

For the public rows, $A=\F_q[S]\cong\F_{q^\ell}$ and
$\dim_{\F_q}A=\ell$.  Write $\Sigma_\xi=I_n\otimes\xi$ for $\xi\in A$ and
use the symmetric scalar expansions
$P_1,\ldots,P_{m_1}\in\Mat_{N_0}(\F_q)$ from
\Cref{eq:scheme-base-public}.  For $\xi,\eta\in A$, define
\begin{equation}
  B_i^{\xi,\eta}(x,y)=x^\transpose\Sigma_\xi P_i\Sigma_\eta y.   \label{eq:base-bilinear}
\end{equation}
Setting $x=y$ gives the quadratic power-pair feature.

\begin{table}[H]
\centering
\caption{Dimensions of the nine public $q=19$ Version~2.3 parameter rows.}
\label{tab:parameters}
\small
\begin{tabular}{@{}c l r r r r@{}}
\toprule
Level & $(v,o,\ell,r)$ & $m_1$ & $M$ & $K$ & $N_0$\\
\midrule
I   & $(28,5,4,4)$  & $5$  & $80$  & $50$ & $132$\\
I   & $(48,16,2,2)$ & $16$ & $64$  & $48$ & $128$\\
I   & $(28,4,4,5)$  & $5$  & $80$  & $50$ & $128$\\
III & $(40,7,4,4)$  & $7$  & $112$ & $70$ & $188$\\
III & $(72,24,2,2)$ & $24$ & $96$  & $72$ & $192$\\
III & $(38,5,4,5)$  & $7$  & $100$ & $70$ & $172$\\
V   & $(50,9,4,4)$  & $9$  & $144$ & $90$ & $236$\\
V   & $(96,32,2,2)$ & $32$ & $128$ & $96$ & $256$\\
V   & $(52,6,4,6)$  & $9$  & $144$ & $90$ & $232$\\
\bottomrule
\end{tabular}
\end{table}

\subsection{The common-column restriction}
Write the stacked signature columns as
\[
  U=[u_0\mid\cdots\mid u_{r-1}],\qquad u_j\in\F_q^{N_0}.
\]
The attacker chooses public multipliers $R_j\in A$ and offsets
$v_j\in\F_q^{N_0}$, with $R_0=1$, and imposes
\begin{equation}
  u_j=\Sigma_{R_j}u+v_j,
  \qquad 0\le j<r.                                  \label{eq:column-relation}
\end{equation}
This relation selects a public subset of possible signatures.  It makes no
assumption about the honest signer and does not change the public verifier.

Before substitution, the verifier feature vector has
$N_{\rm raw}=m_1\ell^2r^2$ labeled coordinates.  Let $z_{\rm ord}(u)$ list
the remaining scalar quadratic forms indexed by ordered pairs of powers of
$S$.  Let $E$ be the public outer map after the fixed permutation that
reconciles the official flattening order with the block-transposed order used
here.  For a fixed relation $R$, let $R_R$ expand $z_{\rm ord}(u)$ into the
raw feature vector, and let $F_{R,v}$ collect the coefficients of the terms
that are linear in $u$.  The substituted verifier is
\begin{equation}
  E\bigl(R_R z_{\rm ord}(u)+F_{R,v}u\bigr)+c_{R,v}=t.
  \label{eq:substituted-verifier}
\end{equation}
Put
\[
 E_R=ER_R,
 \qquad
 \Lambda_{R,v}=EF_{R,v}.
\]
Products of two variable terms give $z_{\rm ord}(u)$.  Products of one
variable term and one offset give the linear term, and products of two offsets
give the constant.

\subsection{How symmetry removes quadratic forms}
For one base form in this restricted scalar map, the ordered label space is
$A\otimes_{\F_q}A$.  The following theorem identifies the pairs that always
give the same polynomial.

\begin{theorem}[Symmetric-square factorization]\label{thm:symmetry-quotient}
Let $A=\F_q[S]$ have dimension $d$, and assume $S^\transpose=S$ and
$P_i^\transpose=P_i$.  For each $i$, the map
\[
  A\otimes_{\F_q}A\longrightarrow \F_q[u]^{\rm hom}_2,
  \qquad
  \xi\otimes\eta\longmapsto u^\transpose\Sigma_\xi P_i\Sigma_\eta u
\]
is invariant under $\xi\otimes\eta\mapsto\eta\otimes\xi$.  It factors
through
\[
  \Sym^2_{\F_q}(A)=
  (A\otimes_{\F_q}A)/\Span\{\xi\otimes\eta-\eta\otimes\xi\}.
\]
Consequently the direct sum over $m_1$ base forms has rank at most
$m_1\binom{d+1}{2}$.  After any public linear map to $M$ outputs, its rank is
at most
\begin{equation}
  \min\left\{M,\;m_1\binom{d+1}{2}\right\}.          \label{eq:symmetry-cap}
\end{equation}
If the minimal polynomial of $S$ has degree $\ell$, then $d=\ell$ and the
ordered feature vector is obtained by duplicating the off-diagonal
coordinates of an unordered feature vector.
\end{theorem}

\begin{proof}
Every element of $A$ is symmetric.  Since the displayed value is a scalar,
\[
 u^\transpose\Sigma_\xi P_i\Sigma_\eta u
 =\bigl(u^\transpose\Sigma_\xi P_i\Sigma_\eta u\bigr)^\transpose
 =u^\transpose\Sigma_\eta P_i\Sigma_\xi u.
\]
Thus every alternating label $\xi\otimes\eta-\eta\otimes\xi$ lies in the
kernel.  The symmetric square has dimension $\binom{d+1}{2}$, and a later
linear map cannot enlarge its image.  When powers of $S$ form a basis, the
coordinate statement is exactly the duplication of each off-diagonal
unordered pair into its two ordered positions.
\end{proof}

The theorem is an upper bound on the common-column scalar image: extra public
dependencies can only reduce that image further.  Every concrete route checks
the exact public rank before solving.  The outer map cannot repair the loss
because, after the relation, it sees only the already-equal polynomial values.
Before the common-column restriction, swapping $(\xi,\eta)$ only transposes
the matrix $U^\transpose\Sigma_\xi P_i\Sigma_\eta U$; it does not make the
matrix entries equal.  The theorem therefore concerns the restricted scalar
verifier, not the unrestricted matrix-valued map.
The factorization itself uses only transposition symmetry.  Odd
characteristic is needed later for the familiar decomposition into symmetric
and alternating tensors and for the polar-form analysis.  The submitted
nonsymmetrized $q=16$ construction is not covered.

\subsection{Mapping the remaining forms to verifier outputs}
Let $z_{\rm sym}(u)$ contain one coordinate for each unordered pair
$0\le a\le b<\ell$ in every base form, and let $D$ duplicate its off-diagonal
coordinates.  By \Cref{thm:symmetry-quotient},
$z_{\rm ord}(u)=Dz_{\rm sym}(u)$.  Define
\begin{equation}
  Q_R=E_RD\in\F_q^{M\times K}.                     \label{eq:QR}
\end{equation}
The columns of $Q_R$ span every output direction reachable by the quotient
quadratics.

\subsection{Separating the quadratic and affine equations}
Choose $W_R$ whose rows form a basis of the left kernel of $Q_R$, and put
\begin{equation}
  C_{R,v}=W_R\Lambda_{R,v}.                          \label{eq:CRv}
\end{equation}
Multiplication by $W_R$ removes every quadratic term and exposes the affine
constraints.  A complementary map $G_R$ keeps the quadratic coordinates.
The two maps together are an invertible change of verifier outputs.

\begin{lemma}[Invertible output split]\label{lem:canonical-residual}
Assume $\rank Q_R=K$.  Choose a public left inverse
$G_R\in\F_q^{K\times M}$ satisfying $G_RQ_R=I_K$.  Then
\[
 \begin{pmatrix}G_R\\W_R\end{pmatrix}\in\F_q^{M\times M}
\]
is invertible.  Writing $b=t-c_{R,v}$, the complete substituted verifier is
equivalent to
\begin{equation}
 C_{R,v}u=W_Rb,
 \qquad
 g_{R,v}(u)=G_Rb,
 \qquad
 g_{R,v}(u)=z_{\rm sym}(u)+G_R\Lambda_{R,v}u.        \label{eq:canonical-split}
\end{equation}
The polar map of $g_{R,v}$ is exactly that of $z_{\rm sym}$.  If $b$ is
uniform in $\F_q^M$, the two right-hand sides in
\eqref{eq:canonical-split} are independent and uniform.
\end{lemma}

\begin{proof}
If a vector $y\in\F_q^M$ is killed by both $G_R$ and $W_R$, then
$W_Ry=0$ implies $y=Q_Ra$ for some $a\in\F_q^K$.  Applying $G_R$ gives
$a=G_RQ_Ra=0$, so $y=0$.  The stacked square matrix is therefore invertible.
Applying its block rows to
$Q_Rz_{\rm sym}(u)+\Lambda_{R,v}u=b$ gives the two displayed systems, and
invertibility gives the converse.  Affine terms have zero polar.  Finally, an
invertible linear transformation sends a uniform vector to a uniform vector
on the product of the two coordinate spaces.
\end{proof}

\begin{proposition}[Exact affine reduction]\label{prop:exact-affine}
Let $\rho_Q=\rank Q_R$, $\lambda=\rank C_{R,v}$, and
$\delta=M-\rho_Q-\lambda$.  Then
\[
  \rank[Q_R\ \Lambda_{R,v}]=\rho_Q+\lambda.
\]
After row reduction, exactly $\delta$ independent affine consistency
conditions remain on the target.  For every target satisfying those
conditions, Gaussian elimination leaves $\rho_Q$ equations of degree at most
two in $N_0-\lambda$ variables, with exactly the same solution set as the
complete substituted verifier.

If $\rho_Q=K$ and $\lambda=M-K$, every target passes the affine consistency
test and the residual system consists of $K$ equations on an affine translate
of $\ker C_{R,v}$.
\end{proposition}

\begin{proof}
The row space of $W_R$ is $\ker Q_R^\transpose$, so
$\ker W_R=(\ker Q_R^\transpose)^\perp=\im Q_R$.  Hence $W_R$ realizes the
quotient map $\F_q^M\to\F_q^M/\im Q_R$.  The image of the columns of
$\Lambda_{R,v}$ in this quotient has dimension
$\rank(W_R\Lambda_{R,v})=\lambda$, proving the first identity.

The projected affine system is $C_{R,v}u=W_Rb$.  Its codomain has dimension
$M-\rho_Q$ and its image has dimension $\lambda$, leaving
$\delta=(M-\rho_Q)-\lambda$ target consistency equations.  Whenever they
hold, parameterize the affine solution space by $N_0-\lambda$ free variables
and substitute into a complementary set of $\rho_Q$ rows.  Every operation is
invertible row reduction or exact affine elimination, so the solution set is
unchanged.  In the full-rank case, $C_{R,v}$ is surjective and $\delta=0$.
\end{proof}

\subsection{Restricting the remaining equations to an affine slice}
The preceding proposition preserves the full solution set inside the
common-column relation.  The attack may now choose a smaller affine space.
This second restriction preserves every verifier equation, but it need not
contain a solution.  Root existence is handled separately in
\Cref{sec:sq-full-domain-spectrum}.

\begin{corollary}[Exact affine restriction]
\label{cor:exact-affine-restriction}
Assume $\rank Q_R=K$ and $\rank C_{R,v}=M-K$.  For every target $t$, choose
one solution $u_\star$ to $C_{R,v}u=W_R(t-c_{R,v})$, an $h$-dimensional
subspace $L\subseteq\ker C_{R,v}$, and an injection
$T:\F_q^h\to L$.  On the affine space $u_\star+L$, the complete substituted
verifier is equivalent to
\begin{equation}
  g_{R,v}(u_\star+Tx)=G_R(t-c_{R,v}),
  \qquad x\in\F_q^h.                               \label{eq:exact-affine-restriction}
\end{equation}
This is a system of $K$ affine quadratic equations in $h$ variables.  Every
root reconstructs, through \eqref{eq:column-relation}, a candidate that
satisfies all equations in the original verifier.
\end{corollary}

\begin{proof}
The rank assumption makes $C_{R,v}$ surjective, so $u_\star$ exists.  The
affine equations in \eqref{eq:canonical-split} hold throughout
$u_\star+L$.  The invertible output change in
\Cref{lem:canonical-residual} leaves exactly the $K$ residual equations in
\eqref{eq:exact-affine-restriction}.  That lemma also proves the claimed
equivalence.
\end{proof}

Taking $h=K$ gives the square restriction used by the earlier solver
analysis.

\begin{corollary}[Exact square restriction]
\label{cor:exact-square-restriction}
Assume
\[
  \rank Q_R=K,
  \qquad
  \rank C_{R,v}=M-K,
  \qquad
  N_0\ge M.
\]
Then, for every target $t$, choose one solution $u_\star$ to
$C_{R,v}u=W_R(t-c_{R,v})$, a $K$-dimensional subspace
$L\subseteq\ker C_{R,v}$, and an isomorphism
$T:\F_q^K\to L$.  On the affine space $u_\star+L$, the complete substituted
verifier is equivalent to the square system
\begin{equation}
  g_{R,v}(u_\star+Tx)=G_R(t-c_{R,v}),
  \qquad x\in\F_q^K.                               \label{eq:exact-square-restriction}
\end{equation}
Every root reconstructs, through \eqref{eq:column-relation}, a candidate that
satisfies the original verifier equations.  Conversely, every verifier
preimage that satisfies both public restrictions gives a root of
\eqref{eq:exact-square-restriction}.  The corollary does not assert that a root
exists or that it passes the signature-format rejection rule.
\end{corollary}

\begin{proof}
The rank assumption on $C_{R,v}$ makes it surjective, so $u_\star$ exists for
every target.  Its kernel has dimension
\[
  N_0-(M-K)=N_0-M+K\ge K,
\]
so it contains a $K$-dimensional subspace $L$.  The affine equations in
\eqref{eq:canonical-split} hold throughout $u_\star+L$.  The invertible output
change in \Cref{lem:canonical-residual} therefore leaves exactly the $K$
quadratic equations in \eqref{eq:exact-square-restriction}.  That lemma also
gives both directions of the claimed equivalence.
\end{proof}

All nine published parameter shapes satisfy $N_0\ge M$ by
\Cref{tab:parameters}.  Thus, for these shapes, the two displayed rank checks
are enough to construct the exact square system.  They are not enough to show
that this system contains an acceptable root.

\begin{theorem}[Exact quotient for all published shapes]
\label{thm:exact-all-nine}
For every published $q=19$ Version~2.3 shape in \Cref{tab:parameters}, the
common-column restriction factors each public base form through
$\Sym^2(A)$.  After the stated public rank checks, the complete verifier is
equivalent to $K=m_1\binom{\ell+1}{2}$ affine quadratic equations on a public
affine space.  For $\ell=4$, this changes sixteen ordered forms to ten per
base form.  For $\ell=2$, it changes four ordered forms to three.  Every
further affine slice preserves all $K$ residual equations, and every returned
candidate is checked by the original verifier.
\end{theorem}

\begin{proof}
Apply \Cref{thm:symmetry-quotient,prop:exact-affine,cor:exact-affine-restriction}
to each row of \Cref{tab:parameters}.
\end{proof}

\paragraph{Operational meaning.}
The attacker computes every matrix in this section from the supplied public
key.  The attacker may test several public relations and offsets before
solving, but may not replace an unfavorable key.  This paper does not prove
how often that search succeeds or measure its cost separately.  Those costs
belong to $C_{\rm aux}$ for the current PXL model.  If no tested relation
passes the two rank checks, the conditional recovery result does not apply to
that key.  Once the checks pass, the steps in this section are exact public
computations and no verifier equation is dropped.

\subsection{A Level-I dimension walkthrough}
For $(v,o,\ell,r)=(28,5,4,4)$, we have $n=33$, $N_0=132$, $M=80$,
$m_1=5$, and $K=50$.  The common-column relation leaves $80$ ordered
quadratic labels, but the symmetry quotient leaves only $50$ distinct
quadratics.  In the full-rank affine case, the other $30$ output directions
become linear equations and the full affine residual domain has dimension
$132-30=102$.  The attacker can therefore choose a $50$-dimensional affine
space and obtain $50$ quadratic equations in $50$ variables.  This is the
exact reduction.  Whether the chosen space contains an acceptable root, and
how much work is needed to find one, are the subjects of the conditional
analysis later in the paper.

\section{Slice preparation and target sampling}
\label{sec:stable-rejection}
\label{sec:sq-preparing}

The exact reduction permits any suitable affine slice inside
$\ker C_{R,v}$.  Recovery needs a slice chosen without reading its fresh
coefficients, signatures that pass the format rule, and uniform targets.  This
section provides those inputs; root existence is addressed in
\Cref{sec:sq-full-domain-spectrum}.

\subsection{An unconditional stable kernel}
Generic offsets can supply the missing affine output directions, but their
linear terms may destroy the block grading used by multihomogeneous solvers.
Choose one vector $w\in\F_q^{N_0}$ and multipliers
$\Gamma_0,\ldots,\Gamma_{r-1}\in A$, and set
\begin{equation}
  v_j=\Sigma_{\Gamma_j}w.                            \label{eq:structured-offsets}
\end{equation}
Define the offset row space
\begin{equation}
 U_w=\Span_{\F_q}\{w^\transpose\Sigma_\xi P_i\Sigma_\eta:
     1\le i\le m_1,\ \xi,\eta\in A\}.              \label{eq:Uw}
\end{equation}
For a basis $1,S,\ldots,S^{d-1}$ of $A$, form the complete right-$A$ orbit
\begin{equation}
 O^F_{R,v}=
 \begin{pmatrix}
  F_{R,v}\\F_{R,v}\Sigma_S\\\vdots\\F_{R,v}\Sigma_{S^{d-1}}
 \end{pmatrix},
 \qquad \rho_F=\rank O^F_{R,v}.                    \label{eq:orbit-matrix}
\end{equation}

\begin{theorem}[Stable-lift kernel]\label{thm:stable-kernel}
Assume $S^\transpose=S$, $P_i^\transpose=P_i$, and $A\cong\F_{q^d}$.  Under
\eqref{eq:structured-offsets}:
\begin{enumerate}[label=(\roman*)]
\item $U_w$ is closed under right multiplication by $A$ and
      $\dim_{\F_q}U_w\le m_1d^2$;
\item every row of $F_{R,v}$ and $O^F_{R,v}$ lies in $U_w$;
\item the row space of $O^F_{R,v}$ is an $A$-subspace, so $d\mid\rho_F$;
\item
\[
 H^F_{R,v}=\ker O^F_{R,v}
\]
is $A$-linear, satisfies $F_{R,v}H^F_{R,v}=0$, and has
\begin{equation}
  \dim_A H^F_{R,v}=n-\rho_F/d\ge n-m_1d.             \label{eq:stable-dim}
\end{equation}
\end{enumerate}
No generic orbit-rank equality is required.  An unexpected dependency only
enlarges the stable space.
\end{theorem}

Right $A$-stability puts the orbit rows in an $A$-linear row space, and
rank--nullity gives \eqref{eq:stable-dim}.  The complete proof is in
\Cref{app:slice-proofs}.

Since $C_{R,v}=W_REF_{R,v}$, the stable kernel lies inside
$\ker C_{R,v}$.  Translation along a subspace $L\subseteq H^F_{R,v}$ changes
the quotient quadratics but leaves every offset-linear feature constant.

\subsection{Reserved-line decoupling and the structural preflight}
Let $e_1,\ldots,e_n$ be the native $A$-coordinate basis in which the
key-dependent symmetric public matrices are expanded.  Fix a vinegar index
$j_0$ and use the coordinate-aligned decomposition
\begin{align}
 W&=Ae_{j_0},&
 V'&=\bigoplus_{\substack{j\ {\rm vinegar}\\j\ne j_0}}Ae_j,&
 O_{\rm pub}&=\bigoplus_{j\ {\rm oil}}Ae_j,\notag\\
 V_{\rm pub}&=W\oplus V'\oplus O_{\rm pub},&
 \dim_AW&=1,&
 \dim_AV'&=v-1,\qquad \dim_AO_{\rm pub}=o.
 \label{eq:public-decomposition}
\end{align}
The actual SHAKE stream is deterministic once the public seed is fixed.  For
probability statements only, the \emph{native-coordinate idealized
XOF-transcript model} replaces the relevant bytes by independent uniform bytes
and reduces the native upper-triangular entries of the key-dependent $P_i$
modulo $19$.  This is not a claim about a fixed SHAKE-generated key.  The model
conditions on $E$, $R$, and the full vector
$\Gamma=(\Gamma_0,\ldots,\Gamma_{r-1})$, and invokes independence before any
$A$-linear basis change.  A reduced byte has maximum atom
\[
 \rho_{\rm byte}=\frac{14}{256},
\]
because nine residues have fourteen preimages and ten have thirteen.

\begin{lemma}[Coordinate-aligned reserved-line decoupling]
\label{lem:reserved-decoupling}
Fix $E,R,\Gamma$, choose $0\ne w\in W$, and use the structured offsets
\eqref{eq:structured-offsets}.  Then $Q_R$ depends only on $E$ and $R$, while
$C_{R,v}$, $O^F_{R,v}$, $H^F_{R,v}$, and every deterministically tie-broken
slice selected from them depend only on native entries having at least one
index in $W$.  In the native-coordinate idealized XOF-transcript model they are
independent of the native $V'\times V'$ entries.
\end{lemma}

Every offset contains one coordinate from $W$, so the construction never reads
a native $V'\times V'$ entry.  The complete proof is in
\Cref{app:slice-proofs}.

\begin{lemma}[Maximum-atom pivot bound]\label{lem:max-atom-pivot}
Let $X_1,\ldots,X_N$ be independent finite-field variables, each with maximum
atom at most $\rho_{\rm byte}$.  If a matrix $B$ has rank $t$, then for every
$c$,
\[
 \Pr[BX=c]\le\rho_{\rm byte}^{\,t}.
\]
\end{lemma}

Conditioning on the nonpivot variables leaves one allowed value for each pivot,
so the probabilities multiply.  The complete proof is in
\Cref{app:slice-proofs}.

Put $\mathcal Q_R=\F_q^M/\im Q_R$.  When $\rank Q_R=K$,
$\dim_{\F_q}\mathcal Q_R=M-K$, and $C_{R,v}$ is a coordinate matrix for the
induced linear map $u\mapsto W_R\Lambda_{R,v}u$ to $\mathcal Q_R$.  Let
$\mathbf X_{W,V'}$ be the vector of native independent public-matrix entries
in the $W\times V'$ blocks.  For
$0\ne\lambda\in\mathcal Q_R^\vee$, define $B_\lambda$ by
\[
 \operatorname{vec}\!\left(
   (\lambda\circ C_{R,v})|_{V'}\right)
 =B_\lambda\mathbf X_{W,V'}+b_\lambda.
\]
Here $B_\lambda$ and $b_\lambda$ are fixed once $E,R,\Gamma,\lambda$ and the
native layout are fixed; they do not depend on the sampled entries in
$\mathbf X_{W,V'}$.
Scaling $\lambda$ scales $B_\lambda$, so its rank depends only on
$[\lambda]$.  Put
\[
 t_{\rm src}
 =\min_{0\ne\lambda\in\mathcal Q_R^\vee}\rank B_\lambda.
\]
For an $A$-line $\mathfrak l\subseteq W\oplus O_{\rm pub}$, choose
$0\ne y_{\mathfrak l}\in\mathfrak l$ and write
\[
 O^F_{R,v}y_{\mathfrak l}
 =D_{\mathfrak l}\mathbf X_{\rm nat}+e_{\mathfrak l}
\]
in the complete vector $\mathbf X_{\rm nat}$ of native independent
public-matrix entries.  The coefficient data
$D_{\mathfrak l},e_{\mathfrak l}$ are deterministic before those entries are
sampled, and the rank of $D_{\mathfrak l}$ is independent of the chosen
generator.  Define
\[
 t_W=\rank D_W,\qquad
 t_{\rm oth}=
 \min_{\substack{\mathfrak l\subseteq W\oplus O_{\rm pub}\\
                   \dim_A\mathfrak l=1,\ \mathfrak l\ne W}}
 \rank D_{\mathfrak l}.
\]
These are exact symbolic coefficient-map ranks determined by the fixed
relation, the complete offset pattern, and the native coordinate layout; they
do not depend on the sampled values of the key coefficients.  Define
\begin{align}
 \delta_{\rm src}
 &=\frac{q^{M-K}-1}{q-1}
   \rho_{\rm byte}^{d(v-1)},                          \label{eq:delta-src}\\
 \delta_{\rm proj}
 &=\rho_{\rm byte}^{m_1\binom{d+1}{2}}+
   \left(\frac{|A|^{1+o}-1}{|A|-1}-1\right)
   \rho_{\rm byte}^{m_1d^2},
 \qquad |A|=q^d.                                      \label{eq:delta-proj}
\end{align}

\begin{theorem}[Conditional coefficient-preflight density]
\label{thm:structural-preflight}
Fix the deterministic outer map $E$, a relation $R$, and the complete vector
$\Gamma$.  Assume the exact deterministic checks
\[
\rank Q_R=K,\qquad
t_{\rm src}\ge d(v-1),\qquad
t_W\ge m_1\binom{d+1}{2},\qquad
t_{\rm oth}\ge m_1d^2
\]
pass.  In the native-coordinate idealized XOF-transcript model,
\[
 \Pr[\rank C_{R,v}<M-K]\le\delta_{\rm src},
\]
and
\[
 \Pr[H^F_{R,v}\cap(W\oplus O_{\rm pub})\ne0]
 \le\delta_{\rm proj}.
\]
Consequently full affine rank and injectivity of
$H^F_{R,v}\to V'$ hold jointly with probability at least
$1-\delta_{\rm src}-\delta_{\rm proj}$.  For the six official $\ell=4$
dimension tuples, direct substitution gives
$\delta_{\rm src}+\delta_{\rm proj}<2^{-209}$.  The coefficient-map ranks are
fixed-layout symbolic checks; the realized quotient, affine, and intersection
ranks are exact public checks on a fixed key.
\end{theorem}

The two bounds apply \Cref{lem:max-atom-pivot} to projective output
directions and to $A$-lines in $W\oplus O_{\rm pub}$, then take union bounds.
The complete proof is in \Cref{app:slice-proofs}.

The concrete check records all $r$ entries of $\Gamma$ and runs the displayed
tests.  This release does not assert that every fixed key passes them.  This
reserved-line preflight is separate from $\mathsf H_{\rm off}$ and
$\mathsf H_{\rm struct}$ and does not establish either PXL route premise.

\subsection{The guaranteed fresh ordinary square}
Let
\[
  D_0=V'\cap\ker C_{R,v}.
\]
The intersection inequality gives
\begin{equation}
  \dim_{\F_{19}}D_0
  \ge 4(v-1)-(M-K).                                  \label{eq:fresh-slice-dim}
\end{equation}
For all six $\ell=4$ rows this lower bound is at least $K$:
\begin{center}
\small
\begin{tabular}{@{}l r r@{}}
\toprule
Parameters & $\dim D_0$ lower bound & $K$\\
\midrule
$(28,5,4,4)$, $(28,4,4,5)$ & $78$ & $50$\\
$(40,7,4,4)$ & $114$ & $70$\\
$(38,5,4,5)$ & $118$ & $70$\\
$(50,9,4,4)$ & $142$ & $90$\\
$(52,6,4,6)$ & $150$ & $90$\\
\bottomrule
\end{tabular}
\end{center}
Consequently the attacker can choose, from cross data alone, an ordinary
$K$-dimensional slice inside $D_0$.  Its internal $V'\times V'$ coefficients
remain fresh for the spectrum theorem.  This supplies an auxiliary square
slice for every $\ell=4$ row, including $(38,5,4,5)$.  The PXL tables use
different dimensions and premises.

\subsection{Passing the block-symmetry rejection rule}
Version~2.3 rejects a square $\ell=4$ signature if any block is symmetric and
rejects an $\ell=2$ signature if more than $\lfloor n/4\rfloor$ blocks are
symmetric~\cite{SNOVA23}.  Rectangular $\ell=4$ blocks are not subject to the
square-block rule.

For a square $\ell=4$ relation $R=(R_0,R_1,R_2,R_3)$, define the local skew
map
\[
 \mathcal S_R:\F_q^4\longrightarrow\bigwedge^2\F_q^4,
 \qquad
 \mathcal S_R(x)_{a,c}=(R_cx)_a-(R_ax)_c.
\]
For any offsets, one signature block is symmetric exactly when
$\mathcal S_R(x)$ equals a fixed offset-dependent vector.  The public relation
$R_j=S^j$ for the official matrix has $\rank\mathcal S_R=4$.  Thus each
block has at most one rejected common-column value, and the accepted density
in the full common-column space is at least
\begin{equation}
  \alpha_{4,n}=(1-19^{-4})^n.                        \label{eq:l4-acceptance}
\end{equation}
The accepted-root theorem later uses only this global density; it assumes no
independence from the MQ system.

For $\ell=2$, the zero-offset relation used in the main theorem gives a block
\[
 \begin{pmatrix}x&0\\y&0\end{pmatrix},
\]
which is symmetric exactly when $y=0$.  The exact accepted density is derived
in \eqref{eq:sq-l2-alpha}.  A deterministic fixed-skew alternative is recorded
in \Cref{app:slice-proofs}; it is not used by the main theorem.

\subsection{Exactly uniform official targets}
The probability theorems require a uniform target over $\F_{19}^M$, but the
official binary-to-field conversion has a small modulo bias.  Model the target
hash/XOF chunks as independent uniform bit strings, separately from the key
expansion above.  The pinned conversion uses each full eight-byte chunk for
fifteen base-$19$ digits and a final partial chunk for the remainder.  If a
uniform $b$-bit chunk $X$ supplies $a$ digits, retain the message--salt input
only when
\begin{equation}
  X<\left\lfloor\frac{2^b}{19^a}\right\rfloor19^a.   \label{eq:target-rejection-sampling}
\end{equation}
Conditioned on this event, $X\bmod19^a$ is uniform.  Applying the test to each
chunk makes the retained official target exactly uniform, and an invertible
output change preserves uniformity.  Exact evaluation over the nine public
Version~2.3 lengths gives a minimum acceptance probability
$0.1524622985\ldots$, and hence at least $0.152462$.
When target consistency requires more than the $128$-bit salt space, the
existential chosen-message attacker varies a message prefix as well as the
salt.  Target generation occurs before nonlinear solving and adds to, rather
than multiplies, the solve cost.


\paragraph{The exact $\ell=2$ zero-offset split.}
For the three official $\ell=r=2$ rows put
\[
 A=\F_{19^2},\qquad m=m_1=o,\qquad M=4m,\qquad K=3m,
\]
and use the public relation
\begin{equation}
 u_0=u,\qquad u_1=0.                                  \label{eq:sq-l2-zero}
\end{equation}

\begin{proposition}[Exact zero-offset target split]
\label{prop:sq-zero-split}
Let $Q_R$ be the quotient-output matrix for \eqref{eq:sq-l2-zero}.  The
route applies to a fixed key only after the exact public test
$\rank Q_R=K$.  After a passing test, choose a left inverse $G_R$ and a
full-row-rank matrix $W_R$ spanning the left kernel of $Q_R$.  The complete
substituted verifier is equivalent to
\begin{equation}
 W_Rt=0,\qquad z_{\rm sym}(u)=G_Rt.                  \label{eq:sq-l2-target-split}
\end{equation}
For a uniform official target, consistency has probability
$19^{-(M-K)}=19^{-m}$.  Conditioned on consistency, the residual target is
uniform in $\F_{19}^K$, so the attacker can filter message and salt inputs
before solving.  Under the stipulated $2^{32}$ target-charge convention in
\eqref{eq:pxl-kernel-cost}, this contributes $2^{32}19^m$ binary-gate units
to the model numerator.  This is an accounting convention, not a measured
target-generation bound.
\end{proposition}

\begin{proof}
Zero offsets remove the affine and constant terms, leaving
$Q_Rz_{\rm sym}(u)=t$.  The invertible output change with block rows
$G_R,W_R$ gives \eqref{eq:sq-l2-target-split}.  It preserves uniformity and
separates the residual target from the consistency coordinates.  Since target
filtering precedes the nonlinear solve, its work adds to the solve ledger.
\end{proof}

In one block, \eqref{eq:sq-l2-zero} has matrix
$\left(\begin{smallmatrix}x&0\\y&0\end{smallmatrix}\right)$ and is symmetric
exactly when $y=0$.  The $n=v+o$ block conditions use disjoint coordinates.
Thus the exact density accepted by the Version~2.3 rule is
\begin{equation}
 \alpha_n=19^{-n}\sum_{j=0}^{\lfloor n/4\rfloor}
          \binom nj18^{n-j}.                          \label{eq:sq-l2-alpha}
\end{equation}

\section{Affine coefficient support for the PXL routes}
\label{sec:pxl-support}

The PXL analysis keeps every residual equation and permits an
overdetermined affine slice.  We first prove full affine coefficient support
for $\ell=2$.  We then give the exact fixed-map criterion that the
$\ell=4$ slice family must satisfy.

\subsection{Transverse slices for \texorpdfstring{$\ell=2$}{ell=2}}
\label{sec:pxl-l2-slice}

For the three official $\ell=r=2$ shapes, put
$A=\F_{19^2}$ and identify $A^{v+o}$ with
$\F_{19}^{2(v+o)}$.  Write
$A^{v+o}=V_{\rm pub}\oplus O_{\rm pub}$, where
$V_{\rm pub}\cong A^v$, and let $\pi_v$ be the projection to
$V_{\rm pub}$.  Under the zero-offset relation
\eqref{eq:sq-l2-zero}, the public rank test in
\Cref{prop:sq-zero-split} leaves $K=3m$ residual equations and an exact
target-consistency test.  For Level I, this is a map
\begin{equation}
 g:A^{64}\longrightarrow\F_{19}^{48}.              \label{eq:pxl-level1-residual}
\end{equation}

An injection $T:\F_{19}^{h}\to A^{v+o}$ defines the direction
$L=\im T$.  For an $\F_{19}$ basis
$\ell_1,\ldots,\ell_h$ of $L$, put
\begin{equation}
 L_A=\Span_A\{\pi_v(\ell_1),\ldots,\pi_v(\ell_h)\}
 \subseteq V_{\rm pub}.                              \label{eq:pxl-extension-span}
\end{equation}
We call $L$ \emph{$A$-independent} if $\dim_A L_A=h$.  We call an affine
coset $a+L$ \emph{transverse} if
\begin{equation}
 \pi_v(a)\notin L_A.                                 \label{eq:pxl-l2-transverse}
\end{equation}
Both conditions are stronger than base-field injectivity alone.

\begin{lemma}[Uniform coset sampling]
\label{lem:pxl-uniform-coset}
Let $U$ be an $\F_q$ space and let $L\subseteq U$ have dimension $h$.
If $a$ is uniform in $U$, then $a+L$ is uniform in $U/L$.  For every
$x\in U$,
\begin{equation}
 \Pr[x\in a+L]=q^{h-\dim U}.                         \label{eq:pxl-point-in-coset}
\end{equation}
If $z$ is independently uniform in $\F_q^K$, then
\begin{equation}
 \Pr[x\in a+L\ \text{and}\ g(x)=z]
 =q^{h-\dim U-K}.                                    \label{eq:pxl-point-root}
\end{equation}
If $U=A^{v+o}$ and $L$ is $A$-independent, then
\begin{equation}
 \Pr[\pi_v(a)\notin L_A]=1-|A|^{h-v}.               \label{eq:pxl-transverse-probability}
\end{equation}
\end{lemma}

\begin{proof}
Every coset has $q^h$ representatives, which proves the first claim and
\eqref{eq:pxl-point-in-coset}.  Independence of $z$ contributes $q^{-K}$ in
\eqref{eq:pxl-point-root}.  The projection $\pi_v(a)$ is uniform in $A^v$,
while $L_A$ has $|A|^h$ elements.  This proves
\eqref{eq:pxl-transverse-probability}.
\end{proof}

\begin{proposition}[Exact transverse affine restriction]
\label{prop:pxl-level1-transverse}
Assume the public rank check in \Cref{prop:sq-zero-split}.  For every
consistent target and every affine coset $a+\im T$, the complete substituted
verifier is equivalent on that coset to
\begin{equation}
 f_{a,z}(x)=g(a+Tx)-z=0,
 \qquad x\in\F_{19}^{h}.                             \label{eq:pxl-restricted-system}
\end{equation}
This is a system of $K$ affine quadratic equations in $h$ variables.  Every
accepted root reconstructs a candidate that satisfies every equation in the
original verifier.
\end{proposition}

\begin{proof}
The target split is an invertible public change of outputs.  The map
$x\mapsto a+Tx$ is an affine injection, so substitution preserves the
solution set in the coset.  The reconstruction uses
\eqref{eq:column-relation}, followed by the format rule and the original
verifier.
\end{proof}

\begin{lemma}[Mixed block support]
\label{lem:pxl-l2-mixed-block}
Let $C:A\times A\to\F_{19}$ be bilinear and put
$C^{\rm sym}(\xi,\zeta)=C(\xi,\zeta)+C(\zeta,\xi)$.  The map
$C\mapsto C^{\rm sym}$ is onto the symmetric bilinear forms on the
two-dimensional $\F_{19}$ space $A$.  Its kernel is the one-dimensional
space of alternating forms.
\end{lemma}

\begin{proof}
Since the characteristic is odd, a symmetric form $S$ has preimage $S/2$.
The kernel consists exactly of alternating forms on $A$.
\end{proof}

\begin{proposition}[$\ell=2$ affine coefficient support]
\label{prop:pxl-l2-affine-support}
Let $h\le v-1$, let $L$ be $A$-independent, and let $a+L$ be transverse.
For one fresh UOV base form, the restriction map from its fresh native
symmetric coefficients to the three descended $\ell=2$ channel polynomials
on $a+L$ is onto
\begin{equation}
  \bigl(\F_{19}[x_1,\ldots,x_h]_{\le2}\bigr)^3.      \label{eq:pxl-l2-affine-surjection}
\end{equation}
Across the $m$ independent base forms, the coefficient map is onto the space
of $K=3m$ ordinary affine quadratic equations in $h$ variables.  In
particular, $h=v-2$ has transversality failure probability
$|A|^{-2}=19^{-4}$.
\end{proposition}

\begin{proof}
Put $b=\pi_v(a)$ and $e_i=\pi_v(\ell_i)$.  Transversality and
$A$-independence give the direct sum
\[
 Ab\oplus Ae_1\oplus\cdots\oplus Ae_h\subseteq V_{\rm pub}.
\]
The diagonal block on $Ae_i$ supplies the three channel coefficients of
$x_i^2$.  The block between $Ae_i$ and $Ae_j$ supplies all three
coefficients of $x_ix_j$ by \Cref{lem:pxl-l2-mixed-block}.  The blocks
between $Ab$ and $Ae_i$ supply the linear coefficients, and the diagonal
block on $Ab$ supplies the constants.  These blocks are independent, so
every triple of affine quadratics has a native preimage.

The block map is defined over $\F_{19}$.  Its surjectivity after scalar
extension therefore descends by faithful flatness.  Taking the direct sum
over the independent base forms gives all $K$ equations.  The final
probability is \eqref{eq:pxl-transverse-probability} with $h=v-2$.
\end{proof}

This proposition proves support, not distribution.  Uniform output follows
from a uniform native source, but the official byte-reduction source is not
uniform on $\F_{19}$.  The proposition does not prove semi-regularity or any
Macaulay rank prediction.

\subsection{Fixed-map support for \texorpdfstring{$\ell=4$}{ell=4}}
\label{sec:pxl-l4-support}

Work over a splitting field $\Omega$ of odd characteristic.  Let the four
Frobenius blocks be $V_a$, with planted oil spaces $O_a\subseteq V_a$, and
fix maps $T_a:X\to V_a$ for $0\le a<4$.  Put
$\overline T_a:X\to V_a/O_a$.  The allowed diagonal coefficient block
consists of symmetric forms on $V_a$ that vanish on $O_a\times O_a$.  The
allowed cross block for $a<b$ consists of bilinear forms on
$V_a\times V_b$ that vanish on $O_a\times O_b$.  Pullback gives a fixed
linear map $\Phi_T$ to four diagonal and six cross quadratic forms on $X$.

\begin{theorem}[Fixed-map channel surjectivity]
\label{thm:pxl-fixed-map-surjectivity}
The map
\[
 \Phi_T:\mathcal U_O\longrightarrow\Sym^2(X^*)^{10}
\]
is onto if and only if every $\overline T_a$ is injective.
\end{theorem}

\begin{proof}
Suppose every $\overline T_a$ is injective.  Choose complements
$V_a=O_a\oplus C_a$.  The projection of $T_a(X)$ to $C_a$ is an isomorphic
copy of $X$.  Any symmetric form on $X$ can therefore be transported to
that copy, extended to $C_a$, and set to zero on blocks involving $O_a$.
This realizes every diagonal channel.

For $a<b$, the two projected copies of $X$ are injective.  Given
$S\in\Sym^2(X^*)$, prescribe the cross bilinear form to be $S/2$ on these
copies, extend it to $C_a\times C_b$, and set blocks involving the oil spaces
to zero.  Its symmetric pullback is $S$.  The ten source blocks are
independent, so their direct sum is onto.

Conversely, if $0\ne x\in\ker\overline T_a$, then $T_ax\in O_a$.
Every allowed diagonal form vanishes at $(T_ax,T_ax)$, so that diagonal
channel cannot realize a quadratic form that is nonzero at $x$.
\end{proof}

For both Level I $\ell=4$ shapes, $\dim(V_a/O_a)=28$.  A fixed PXL main
space of dimension $31$ therefore cannot satisfy this criterion.  Full
support through this theorem requires a main space of dimension at most
$28$, together with a realizable slice for which all four quotient maps are
injective.  If the spaces, maps, and coefficient conditions arise by scalar
extension from $\F_{19}$, surjectivity descends by faithful flatness.  An
eigenblock decomposition over $\Omega$ alone does not supply that descent or
a base-field realization.

The official change of variables supplies the required realization.  Write
the public coordinate space as
\[
 E=U_{\rm vin}\oplus O_{\rm pub},
 \qquad U_{\rm vin}=A^v\oplus0,
 \qquad O_{\rm pub}=0\oplus A^o,
\]
and put $R=T^{12}$.  The hidden oil space in public coordinates is the graph
\begin{equation}
 G=\{(Rz,z):z\in A^o\}.                              \label{eq:pxl-hidden-oil-graph}
\end{equation}

\begin{lemma}[Official graph complement]
\label{lem:pxl-graph-complement}
The public vinegar space is a complement to the hidden oil graph:
\begin{equation}
 E=U_{\rm vin}\oplus G.                              \label{eq:pxl-graph-complement}
\end{equation}
After scalar extension to a splitting field, let $E_a,U_a,G_a$ denote the
four Frobenius components.  For every $a$, the quotient map
\begin{equation}
 U_a\longrightarrow E_a/G_a                         \label{eq:pxl-split-quotient-iso}
\end{equation}
is an isomorphism.
\end{lemma}

\begin{proof}
Every $(x,z)\in E$ has the unique decomposition
\[
 (x,z)=(x-Rz,0)+(Rz,z).
\]
This proves \eqref{eq:pxl-graph-complement}.  Scalar extension and projection
to each primitive idempotent preserve direct sums.  Thus
$E_a=U_a\oplus G_a$, which is equivalent to
\eqref{eq:pxl-split-quotient-iso}.
\end{proof}

Choose $d\le v$ public vinegar coordinate vectors that are $A$-independent,
and let $X$ be their $\F_{19}$ span.  Each split projection of $X$ maps
injectively to $E_a/G_a$ by \Cref{lem:pxl-graph-complement}.  Therefore
\Cref{thm:pxl-fixed-map-surjectivity} applies simultaneously to all four
Frobenius components.  This proves that the official realization family meets
the fixed-map good locus whenever the PXL main dimension is at most $v$.

For a full PXL slice, choose a direct sum
\begin{equation}
 L=X\oplus Z,
 \qquad \dim X=d,
 \qquad \dim Z=k,                                   \label{eq:pxl-xz-slice}
\end{equation}
and a public coordinate map $\Gamma:L\to\F_{19}^k$ that kills $X$ and is an
isomorphism on $Z$.  Fixing the $\Gamma$ coordinates leaves an affine
translate of $X$.  Its leading quadratic forms are the pullbacks analyzed by
\Cref{thm:pxl-fixed-map-surjectivity}.  The constrained ledger below chooses
$d=h-k=v$ for every row.

Zero offsets put $C_{R,0}=0$, so every such $L$ lies in the residual affine
kernel.  To avoid the resulting target-consistency cost, the supported route
uses the following separate offset premise.

\begin{definition}[Oil-supported offset premise $\mathsf H_{\rm off}$]
\label{def:pxl-offset}
For the selected public relation, there are publicly specified oil-supported
offsets for which the exact checks
\begin{equation}
 \rank(C_{R,v}|_{U_{\rm vin}})=\lambda,
 \qquad
 A\,\ker(C_{R,v}|_{U_{\rm vin}})=U_{\rm vin}         \label{eq:pxl-offset-checks}
\end{equation}
hold, and the construction in \eqref{eq:pxl-xz-slice} can be performed inside
$\ker C_{R,v}$, where $\lambda=M-K$ is the full-rank value used by the
ledger.  Here $A\,W$ denotes the $A$ span of an $\F_{19}$ subspace
$W$.  The matrices and both displayed conditions are exact public checks.
The premise makes no claim about how often official keys or sampled offsets
pass them.
\end{definition}

Under $\mathsf H_{\rm off}$, the graph complement proves homogeneous channel
support on the PXL main variables.  It does not prove full affine coefficient
support.  Semi-regularity, Macaulay ranks, and extraction remain conditional
on $\mathsf H_{\rm route}$.  The frequency of the offset premise under the
official key distribution is unproved.

\section{Root existence on an affine slice}
\label{sec:sq-full-domain-spectrum}

The preceding sections produce affine MQ systems but do not show that a chosen
slice has a root with the required format and a full-column-rank Jacobian.
This section derives that probability from an exact collision identity and
bounds the required derivative ranks for the slices used later.

\subsection{Polar forms, rank spectra, and collisions}
For a scalar quadratic polynomial $f$, define the polar form
\[
  \beta_f(x,y)=f(x+y)-f(x)-f(y)+f(0).
\]
In odd characteristic this is bilinear.  For a vector-valued map $g$, a
nonzero $b\in\F_q^K$ selects the scalar quadratic $b\cdot g$.  Its polar rank
measures how many independent directions the quadratic part couples.

For a fixed slice direction $y$, the vector-valued derivative
\begin{equation}
  T_y(x)=g(x+y)-g(x)-g(y)+g(0)
  \label{eq:sq-Ty}
\end{equation}
is linear in $x$.  The value $q^{-\rank T_y}$ is exactly the fraction of its
codomain-dual directions that annihilate $T_y$, after the appropriate
normalization.  Summing these quantities over nonzero $y$ controls the second
factorial moment of the number of roots in a random target--coset pair.  The
same spectrum controls the density of points where the restricted Jacobian
loses rank.

\subsection{The exact first and second moments}

\begin{theorem}[Root-collision identity]\label{thm:collision-identity}
Let $U$ be an $\F_q$-space, let $C:U\to\F_q^c$ be surjective, let
$L\subseteq\ker C$ have dimension $h$, and let
$g:U\to\F_q^K$ have degree at most two.  Sample independent uniform
$d\in\F_q^c$ and $z\in\F_q^K$, then choose a uniform affine coset $\mathcal A$
of $L$ inside $C^{-1}(d)$.  Put
\[
 Z=|\{x\in\mathcal A:g(x)=z\}|,
 \qquad \mu=q^{h-K}.
\]
For $0\ne y\in L$, define
\[
 T_y(x)=g(x+y)-g(x)-g(y)+g(0)
\]
and let $\kappa(y)$ be one if
$-(g(y)-g(0))\in\im T_y$ and zero otherwise.  Then
\begin{align}
 \mathbb EZ&=\mu,                                      \label{eq:moment-first}\\
 \mathbb E[Z(Z-1)]
 &=\mu\sum_{0\ne y\in L}\kappa(y)q^{-\rank T_y}
 \le \mu\sum_{0\ne y\in L}q^{-\rank T_y}.          \label{eq:moment-second}
\end{align}
The domain of $T_y$ is the full space $U$ because the experiment averages over
the affine right-hand side as well as the coset.
\end{theorem}

\begin{proof}
Sampling $d$ and then a uniform $L$-coset in $C^{-1}(d)$ is exactly uniform
over all cosets of $L$ in $U$.  Every point of $U$ therefore lies in the
sampled coset with probability $q^{h-\dim U}$, and its image matches the
independent target with probability $q^{-K}$.  Summing over all points gives
\eqref{eq:moment-first}.

For an ordered pair of distinct roots in one coset, write the second point as
$x+y$ with $0\ne y\in L$.  The collision condition is
\[
 g(x+y)=g(x)
 \quad\Longleftrightarrow\quad
 T_y(x)=-(g(y)-g(0)).
\]
It has $\kappa(y)q^{\dim U-\rank T_y}$ solutions in $x$.  Divide the total
number of ordered collisions by the number
$q^{\dim U-h}q^K$ of target--coset pairs to obtain
\eqref{eq:moment-second}.
\end{proof}

\subsection{Accepted roots with full column rank}

Retain $U,C,g,L$ and the trial from \Cref{thm:collision-identity}, and write
$g=q_2+q_1+q_0$ by homogeneous degree.  For the full-domain derivative $T_y$
from \eqref{eq:sq-Ty}, define
\begin{equation}
 \bar\eta_L^U=\sum_{0\ne y\in L}q^{-\rank T_y}. \label{eq:sq-etaU}
\end{equation}
The superscript records that the trial averages over the affine right-hand side
and the $L$-coset.  The derivative therefore has domain $U$, not $\ker C$.
Restricting it to $\ker C$ gives a different, generally larger quantity used
for fixed-fiber uniqueness, not for this experiment.

Let $X\subseteq U$ be the assignments accepted by the original verifier and
put $\alpha=|X|/|U|$.  A trial samples independent uniform
$d\in\F_q^c$ and $z\in\F_q^K$, and then a uniform affine coset of $L$ inside
$C^{-1}(d)$.

\begin{theorem}[Accepted root with full column rank]
\label{thm:sq-accepted-root}
Assume $h\le K$, $\bar\eta_L^U\le\eta$, and
$\alpha>\eta/(q-1)$.  A trial contains an accepted root $x$ satisfying
\[
 \rank D_Lg(x)=h,
 \qquad
 D_Lg(x)[y]=T_y(x)+q_1(y),
\]
with probability at least
\begin{equation}
 p_{\rm ns}
 =q^{h-K}\frac{a^2}{a+\eta},
 \qquad a=\alpha-\frac{\eta}{q-1}.
 \label{eq:sq-pns}
\end{equation}
No independence among acceptance, singularity, the target, and the quadratic
map is assumed.
\end{theorem}

\begin{proof}
Let $Y$ count accepted roots in the sampled target--coset pair at which
$D_Lg$ has rank $h$, and let $Z$ count all roots.  For a projective direction
$[y]\in\mathbb P(L)(\F_q)$, the condition $D_Lg(x)[y]=0$ is an affine system
in $x$ whose linear part is $T_y$.  It therefore holds on at most a
$q^{-\rank T_y}$ fraction of $U$.  Union-bound over projective directions to
obtain
\[
 \frac{|\{x:\rank D_Lg(x)<h\}|}{|U|}
 \le \frac1{q-1}\sum_{0\ne y\in L}q^{-\rank T_y}
 \le \frac\eta{q-1}.
\]
Every point of $U$ is selected by the random fiber--coset with probability
$q^{h-\dim U}$ and matches the independent target with probability $q^{-K}$.
Put $\mu=q^{h-K}$ and $a=\alpha-\eta/(q-1)$.  Hence
$\mathbb EY\ge\mu a$.
The ordered-pair collision identity gives
$\mathbb E[Z(Z-1)]\le\mu\eta$, and
$Y(Y-1)\le Z(Z-1)$.  Therefore
\[
 \mathbb E[Y^2]
 =\mathbb EY+\mathbb E[Y(Y-1)]
 \le \mathbb EY+\mu\eta.
\]
The function $x\mapsto x^2/(x+\mu\eta)$ is increasing for $x\ge0$.
Cauchy--Schwarz therefore gives
\[
 \Pr[Y>0]\ge
 \frac{(\mathbb EY)^2}{\mathbb EY+\mu\eta}
 \ge \mu\frac{a^2}{a+\eta},
\]
which proves \eqref{eq:sq-pns}.  The denominator is strictly smaller than
$1+\eta$ whenever $a<1$.
\end{proof}

The following lemma bounds this spectrum in the idealized transcript model.
Let $\rho_{\rm byte}=14/256$ be the maximum atom of one public-expansion byte
reduced modulo $19$.

\begin{lemma}[Full-domain fresh-spectrum theorem]
\label{lem:sq-full-domain-spectrum}
Assume $q$ is odd and $A=\F_{q^d}$.  Work in the native-coordinate idealized
XOF-transcript model of \Cref{sec:stable-rejection}: the scalar entries in the
fresh symmetric-matrix region are independent reductions of independent uniform
bytes, each having maximum atom $\rho_{\rm byte}=14/256$, while all disjoint
coefficient data are conditioned upon.  Let $L$ be fixed from coefficient data
disjoint from that fresh region.  Suppose every $0\ne y\in L$
has nonzero projection $\bar y$ to a fresh $A$-space $D$, and every nonzero
residual-output combination contains a nonzero coefficient
$c\in\operatorname{Sym}_{\F_q}^2(A)$ in at least one fresh public base form.
If
\[
 \dim_{\F_q}D=d+t,
\]
then for all $0\ne y\in L$ and $0\ne b\in\F_q^K$,
\begin{equation}
 \Pr[b\in\ker T_y^{\mathsf T}]\le\rho_{\rm byte}^t.             \label{eq:sq-pointwise}
\end{equation}
Consequently
\begin{align}
 \mathbb E[\bar\eta_L^U]
 &\le (q^h-1)q^{-K}
 +(q^h-1)(1-q^{-K})\rho_{\rm byte}^t,                            \label{eq:sq-eta-expect}\\
 \Pr\!\left[
  \bar\eta_L^U>(q^h-1)q^{-K}+\varepsilon
 \right]
 &\le
 \frac{(q^h-1)(1-q^{-K})\rho_{\rm byte}^t}{\varepsilon}.
                                                                    \label{eq:sq-eta-tail}
\end{align}
For the reserved-line construction one may take
$D=V'\cong A^{v-1}$ and $t=d(v-2)$.  For the zero-offset construction on the
full public-vinegar summand one may take $D\cong A^v$ and $t=d(v-1)$.
\end{lemma}

For fixed $y\ne0$ and $b\ne0$, the fresh coefficient map has rank at least
$t$.  The maximum-atom bound gives \eqref{eq:sq-pointwise}; expectation and
Markov's inequality give \eqref{eq:sq-eta-expect} and
\eqref{eq:sq-eta-tail}.  The complete proof is in
\Cref{app:slice-proofs}.

For the rest of the paper, use the threshold
\begin{equation}
 \boxed{
 \eta_{128}(h,K)=(q^h-1)q^{-K}+2^{-128}.}
 \label{eq:sq-eta128}
\end{equation}
This allows $2^{-128}$ above the exact full-rank baseline; it does not assume
perfect rank.

\begin{corollary}[Numerical spectrum envelope]
\label{cor:sq-spectrum-tails}
For any invoked PXL slice, let $h$ be its base-field dimension and let $t$ be
the number of fresh scalar pivots supplied by
\Cref{lem:sq-full-domain-spectrum}.  Substituting these two values and
$\varepsilon=2^{-128}$ into \eqref{eq:sq-eta-tail} gives the corresponding
base-two failure exponent.  The PXL ledgers record $h$, while the selected
fresh region determines $t$.  A fixed key can instead publish its exact
projective rank histogram, which certifies \eqref{eq:sq-eta128} without the
idealized XOF-transcript model.
\end{corollary}

\section{Conditional PXL recovery}
\label{sec:pxl-l2-recovery}

The exact quotient applies to all nine parameter shapes.  This section gives
the audited numerical solver model for both extension degrees.  The
$\ell=2$ route has exact affine coefficient support.  For $\ell=4$, we give
one route supported by the graph-complement theorem and a separate cheaper
route with a named structured-system premise.

\subsection{The \texorpdfstring{$\ell=2$}{ell=2} route}

For one $\ell=2$ shape, set $h=v-2$ and $K=3m$.  The exact accepted density
on $A^{v+o}$ is
\begin{equation}
 \alpha_{v+o}=19^{-(v+o)}
 \sum_{j=0}^{\lfloor(v+o)/4\rfloor}
 \binom{v+o}{j}18^{v+o-j}.                          \label{eq:pxl-alpha}
\end{equation}
Put $\eta=\eta_{128}(h,K)$ and define
\begin{align}
 a&=\alpha_{v+o}-\frac{\eta}{18},\notag\\
 p_{\rm reg}&=19^{h-K}\frac{a^2}{a+\eta},\notag\\
 p_{\rm dir}&=\max\{p_{\rm reg}-19^{-4},0\}.       \label{eq:pxl-pdir}
\end{align}
The final subtraction is a union bound with the exact nontransverse
probability.  It does not assume independence and is not claimed to be the
exact usable trial probability.

The solver analysis uses two named premises.  The ordinary MQ hypothesis
$\mathsf H_{\rm PXL}$ packages the affine semi-regularity, projected
semi-regularity, and Macaulay-rank assumptions from PXL Section~4.2
~\cite{FurueKudo2024PXL}.  The manuscript adds an instancewise premise
$\mathsf H_{\rm route}$ stating that the actual restricted matrices follow
those predictions and that the added extraction procedure returns an
accepted base-field root when one exists.  The definitions and the complete
five-component cost expression are in \Cref{app:pxl-model}.

\begin{theorem}[Conditional $\ell=2$ PXL estimates]
\label{thm:pxl-l2-estimates}
Consider one of the three published $\ell=2$ shapes.  Assume the exact public
rank check, the spectrum inequality $\bar\eta_L^U\le\eta_{128}(h,K)$ for
every invoked slice, and $\mathsf H_{\rm route}$ for every invoked restricted
system.  Normalize each $\F_{19}$ operation to the $150$ Boolean gates in the
released scalar multiplication circuit.  This is a multiplication-based
normalization, not a bound for every operation used by PXL.
The conservative sum of the five components from PXL Section~4.2, target
generation, and retries using \eqref{eq:pxl-pdir} gives the following
unit-leading gate exponents.

\begin{center}
\small
\setlength{\tabcolsep}{4pt}
\begin{tabular}{@{}c l r r r r r r@{}}
\toprule
$\omega$ & Level and parameters & $K$ & $h$ & $k$ & $D$ & $\log_2 W_{\rm est}$ & headroom\\
\midrule
$2.81$ & I, $(48,16,2,2)$ & 48 & 46 & 9  & 12 & 128.833 & 14.167\\
       & III, $(72,24,2,2)$ & 72 & 70 & 13 & 17 & 187.030 & 19.970\\
       & V, $(96,32,2,2)$ & 96 & 94 & 15 & 23 & 245.042 & 26.958\\
\midrule
$3$    & I, $(48,16,2,2)$ & 48 & 46 & 9  & 12 & 133.667 & 9.333\\
       & III, $(72,24,2,2)$ & 72 & 70 & 14 & 16 & 194.626 & 12.374\\
       & V, $(96,32,2,2)$ & 96 & 94 & 19 & 20 & 255.550 & 16.450\\
\bottomrule
\end{tabular}
\end{center}

The headroom uses the NIST classical reference exponents $143$, $207$, and
$272$~\cite{NISTSecurityCriteria2026}.  The estimate omits setup, extraction,
control, storage traffic, and final verification.  If $C_{\rm aux}$ charges
those operations per trial, \Cref{eq:pxl-augmented-work} gives the augmented
expression.  Every returned candidate is checked against all residual
equations, the signature format rule, and the original verifier.  The route
therefore never returns an invalid signature, while its completeness and cost
remain conditional on $\mathsf H_{\rm route}$.
\end{theorem}

The first theorem row uses $D=12$ and residual matrix size
$A=99{,}219{,}059$.  The unit-leading dense symbolic storage proxy is
$2^{93.337}$ bits.  This counts one dense symbolic array at five bits per
base-field entry.  It omits representation overhead, temporary arrays, and
storage traffic, and is not a peak-memory bound.  The corresponding residual
matrix sizes and storage proxies for Levels III and V are recorded in
\Cref{app:pxl-model}.

\subsection{The support-constrained \texorpdfstring{$\ell=4$}{ell=4} route}
\label{sec:pxl-l4-supported}

For either $\ell=4$ route, put
\begin{equation}
 \alpha_4=
 \begin{cases}
  (1-19^{-4})^{v+o},&r=4,\\
  1,&r>4,
 \end{cases}
 \qquad
 \eta=\eta_{128}(h,K),
 \qquad
 p_4=19^{h-K}\frac{(\alpha_4-\eta/18)^2}
                         {\alpha_4-\eta/18+\eta}.   \label{eq:pxl-l4-pdir}
\end{equation}
This is the accepted-root lower bound from \Cref{thm:sq-accepted-root}.
Every $\ell=4$ row conditions on the displayed spectrum inequality for its
slice and divides by $p_4$.  The rectangular shapes have no square-block
rejection rule, so their accepted density is one.

The supported route imposes $h-k\le v$.  It assumes the exact public checks
in $\mathsf H_{\rm off}$, the spectrum inequality
$\bar\eta_L^U\le\eta_{128}(h,K)$, and the instancewise PXL premise
$\mathsf H_{\rm route}$.  The graph complement and fixed-map theorem give
homogeneous channel support on the main variables.  They do not give full
affine coefficient support.  Charging the same five PXL components, target
generation, and certified retries gives the following unit-leading gate
exponents.

\begin{center}
\scriptsize
\setlength{\tabcolsep}{3pt}
\begin{tabular}{@{}c l r r r r r r r@{}}
\toprule
$\omega$ & parameters & $K$ & $h$ & $k$ & $h-k$ & $D$
& $\log_2 W_{\rm est}$ & $\log_2 M_{\rm sym}^{\rm proxy}$\\
\midrule
$2.81$ & $(28,5,4,4)$  & 50 & 35 & 7 & 28 & 7  & 145.868 & 59.429\\
       & $(28,4,4,5)$  & 50 & 35 & 7 & 28 & 7  & 145.867 & 59.429\\
       & $(40,7,4,4)$  & 70 & 48 & 8 & 40 & 9  & 197.338 & 78.763\\
       & $(38,5,4,5)$  & 70 & 46 & 8 & 38 & 8  & 201.265 & 71.892\\
       & $(50,9,4,4)$  & 90 & 59 & 9 & 50 & 10 & 252.203 & 91.085\\
       & $(52,6,4,6)$  & 90 & 61 & 9 & 52 & 10 & 249.539 & 92.109\\
\midrule
$3$    & $(28,5,4,4)$  & 50 & 35 & 7 & 28 & 7  & 148.400 & 59.429\\
       & $(28,4,4,5)$  & 50 & 35 & 7 & 28 & 7  & 148.400 & 59.429\\
       & $(40,7,4,4)$  & 70 & 49 & 9 & 40 & 9  & 201.379 & 79.763\\
       & $(38,5,4,5)$  & 70 & 46 & 8 & 38 & 8  & 205.022 & 71.892\\
       & $(50,9,4,4)$  & 90 & 60 & 10 & 50 & 10 & 257.209 & 92.085\\
       & $(52,6,4,6)$  & 90 & 61 & 9 & 52 & 10 & 254.991 & 92.109\\
\bottomrule
\end{tabular}
\end{center}

For $\omega=2.81$, the supported route exceeds the Level I reference by
about $2.87$ bits.  It remains below the Level III and Level V reference
exponents.  This is the cost of enforcing the fixed-map support criterion,
not a weakness in the exact $16$-to-$10$ quotient.  The official frequency
of $\mathsf H_{\rm off}$ remains unproved.

\subsection{The cheaper structured \texorpdfstring{$\ell=4$}{ell=4} route}
\label{sec:pxl-l4-structured}

The cheaper route leaves more than $v$ main variables at Level I.  By
\Cref{thm:pxl-fixed-map-surjectivity}, it cannot have full ordinary coefficient
support for any fixed realization.  Full support is sufficient for the usual
random MQ explanation, but it is not necessary for a structured system to
have the predicted Hilbert series and Macaulay ranks.  We therefore state the
needed premise directly.

\begin{definition}[Structured PXL premise $\mathsf H_{\rm struct}$]
\label{def:pxl-structured}
For every selected $\ell=4$ restricted system, the structured Frobenius
channel family has the Hilbert series, solving degree, Macaulay ranks, and
root-extraction behavior used by the five-component PXL model at the stated
parameters.  This premise does not assert ordinary coefficient support or a
uniform MQ distribution.
\end{definition}

The premise is specific to this manuscript.  Its form is comparable to the
residual genericity assumptions used when Hashimoto-style and generalized
partitioning methods replace a large underdetermined system by smaller MQ
subproblems~\cite{Hashimoto2023,AsanumaEtAl2026}.  The known support
obstruction gives it a greater evidential burden than the ordinary PXL model.
It should therefore be tested directly on the structured family.

Under $\mathsf H_{\rm struct}$ and the same spectrum condition used in
\eqref{eq:pxl-l4-pdir}, the audited unit-leading gate exponents are as follows.
The table does not infer ordinary MQ support from these parameters.

\begin{center}
\scriptsize
\setlength{\tabcolsep}{4pt}
\begin{tabular}{@{}c l r r r r r r@{}}
\toprule
$\omega$ & parameters & $K$ & $h$ & $k$ & $D$
& $\log_2 W_{\rm est}$ & $\log_2 M_{\rm sym}^{\rm proxy}$\\
\midrule
$2.81$ & $(28,5,4,4)$  & 50 & 48 & 10 & 12 & 134.038 & 95.266\\
       & $(28,4,4,5)$  & 50 & 48 & 10 & 12 & 134.038 & 95.266\\
       & $(40,7,4,4)$  & 70 & 68 & 12 & 17 & 182.007 & 135.896\\
       & $(38,5,4,5)$  & 70 & 68 & 12 & 17 & 182.006 & 135.896\\
       & $(50,9,4,4)$  & 90 & 88 & 15 & 21 & 230.445 & 172.076\\
       & $(52,6,4,6)$  & 90 & 88 & 15 & 21 & 230.444 & 172.076\\
\midrule
$3$    & $(28,5,4,4)$  & 50 & 48 & 10 & 12 & 138.912 & 95.266\\
       & $(28,4,4,5)$  & 50 & 48 & 10 & 12 & 138.912 & 95.266\\
       & $(40,7,4,4)$  & 70 & 68 & 13 & 16 & 189.638 & 131.353\\
       & $(38,5,4,5)$  & 70 & 68 & 13 & 16 & 189.637 & 131.353\\
       & $(50,9,4,4)$  & 90 & 88 & 18 & 19 & 240.174 & 162.956\\
       & $(52,6,4,6)$  & 90 & 88 & 18 & 19 & 240.174 & 162.956\\
\bottomrule
\end{tabular}
\end{center}

Every row is reproduced by the final independent ledger checker.  The storage
quantity is the exact five-bit dense-array proxy defined in
\eqref{eq:pxl-storage}; it is not total state or peak memory.

\begin{theorem}[Conditional structured PXL estimates for all nine shapes]
\label{thm:pxl-all-nine-structured}
Assume the exact public rank and spectrum checks for every invoked route.
Assume $\mathsf H_{\rm route}$ for the three $\ell=2$ shapes and
$\mathsf H_{\rm struct}$ for the six $\ell=4$ shapes.  Under the
five-component PXL model with $\omega=2.81$, every displayed row lies below
the NIST reference exponent for its level.  Every returned candidate is
checked against the complete residual system, the signature format rule, and
the original verifier.  Thus validity is exact, while completeness and the
unit-leading gate exponents are conditional on the stated solver premises.
\end{theorem}

\section{Implementation, evidence, and reproducibility}
\label{sec:implementation}

The companion artifact provides formula checkers, one official-key reduction
test, and one toy end-to-end forgery test.  \Cref{tab:evidence-hierarchy}
states what each establishes.  None executes a production-size nonlinear
solver.


Both correspondence tests use a KAT-anchored Python transcription.  The
official-key test compares the substituted map with a separate literal
evaluator but stops before the nonlinear solve.  The toy forger reconstructs
and serializes a signature and checks the rejection rule, but its forger and
verifier share scheme helpers and do not call the official C verifier.  It is
therefore composition evidence rather than an independent implementation.

The release does not test $\mathsf H_{\rm PXL}$,
$\mathsf H_{\rm route}$, or $\mathsf H_{\rm struct}$, and it does not measure
the official frequency of $\mathsf H_{\rm off}$ or the omitted auxiliary work.
The numerical model charges the five PXL components, target generation, and
certified retries.  Control, data movement, indexing, extraction, final
verification, and peak state remain outside the unit-leading gate estimate
(\Cref{app:pxl-model}).

A production fixed-key record would need the exact preflights, a certified
slice spectrum, the complete solver and filtering transcript, the reconstructed
signature, and the original verifier transcript.  Only an instrumented run
that records the omitted work could measure total work and memory.

\begin{table}[t]
\centering
\caption{Evidence shipped with the release.  ``Included'' means that the
named script and its pinned input are present, not that the headline recovery
cost has been measured.}
\label{tab:evidence-hierarchy}
\small
\renewcommand{\arraystretch}{1.14}
\begin{tabularx}{0.98\textwidth}{@{}>{\raggedright\arraybackslash}p{0.19\textwidth}>{\raggedright\arraybackslash}p{0.16\textwidth}X@{}}
\toprule
Evidence & Availability & What it establishes\\
\midrule
Formula and ledger checks & Included & The independent PXL ledger generator and checker reproduce all 30 rows, comprising six $\ell=2$, twelve support-constrained $\ell=4$, and twelve structured $\ell=4$ rows, including matrix sizes and the five-bit storage proxy.\\
Scalar circuit netlist & Included & The emitted $150$-gate $\F_{19}$ multiplication netlist is correct on all $361$ valid inputs.  The extension-field figures are recurrence counts, not emitted netlists.\\
Official KAT verification and reduction & Included for one Level-I $(28,5,4,4)$ key & A Python transcription parses the compressed public key, accepts the supplied signature, rejects an altered message, and obtains quotient/affine ranks $50/30$.  The resulting $50$-equation system in $102$ variables is left unsolved.\\
Reduced zero-offset forgery & Included at unofficial $(2,1,2,2)$ & The complete restricted map is interpolated; its image equals the explicit rank-three quotient image, leaving one target-consistency coordinate before public slice search.  The main case, eight committed fresh-key regressions, and a separate 24-forgery/100-rank stress check pass in the shared KAT-anchored Python transcription.  This is toy composition evidence, not independent C-verifier or production-cost evidence.\\
All-shape official-key preflight census & Not included & No empirical pass probability for the reserved-line spectrum checks or $\mathsf H_{\rm off}$.\\
Production-size solver or forgery & Not included or claimed & No empirical evidence for $\mathsf H_{\rm PXL}$, $\mathsf H_{\rm route}$, $\mathsf H_{\rm struct}$, peak memory, or feasibility at a published shape.\\
\bottomrule
\end{tabularx}
\end{table}

\FloatBarrier

\section{Necessary dimension floors within the common-column model}
\label{sec:upgrade-repair}

We give two necessary output floors for redesigns that retain the
common-column architecture.  The first removes a square slice from every
target fiber.  The stronger second floor removes the corresponding
zero-offset chosen-message slice.  Neither floor is sufficient for security.

\begin{proposition}[Dimensionally square common-column slice bound]
\label{prop:upgrade-fresh-square-floor}
Let $A=\F_{q^d}$ and retain the public vinegar summand
$V'$ of base-field dimension $d(v-1)$.  Suppose an exact affine reduction has
$M'$ outputs, $K'$ quotient quadratics, and full affine rank $M'-K'$.  Then
\[
 \dim_{\F_q}(V'\cap\ker C_{R,v})
 \ge d(v-1)-(M'-K').
\]
If $M'\le d(v-1)$, dimensions alone leave an ordinary $K'$-dimensional slice
carrying all $K'$ residual equations.  Equality of equation and variable
counts does not establish a root, zero-dimensionality, nonsingularity, or an
efficient solve.  Eliminating this dimensional guarantee requires
$M'>d(v-1)$.  Under the coupling $M'=dx'$, this forces $x'\ge v$ and
$M'\ge dv$.
\end{proposition}

\begin{proof}
Use
\[
 \dim(V'\cap\ker C_{R,v})
 \ge\dim V'-\rank C_{R,v}
 =d(v-1)-(M'-K').
\]
The right side is at least $K'$ exactly when $M'\le d(v-1)$.  Negate and
round to the next multiple of $d$.
\end{proof}

Let $K_{\rm cc,ord}=\min\{M,m_1d^2\}$ be the pre-symmetrization ordered-label
cap within the same common-column model.  Retaining symmetric forms requires
$m_1'\ge\lceil K_{\rm cc,ord}/\binom{d+1}{2}\rceil$ and
$M'\ge K_{\rm cc,ord}$.  Under
$M'=dx'$ and $m_1'=\lceil x'/d\rceil$, all three necessities combine as
follows.

\begin{theorem}[Combined necessary common-column output floor]
\label{thm:upgrade-combined-floor}
Any redesign that retains symmetric public forms, the public vinegar count,
the public-vinegar affine geometry above, full affine rank $M'-K'$, and the
Version~2.3 output/base-form coupling, while both restoring
$K_{\rm cc,ord}$ and eliminating the dimensionally guaranteed fresh slice,
must satisfy
\begin{equation}
 x'\ge\max\left\{
 v,
  \left\lceil\frac{K_{\rm cc,ord}}d\right\rceil,
 d\left(
  \left\lceil\frac{K_{\rm cc,ord}}{\binom{d+1}{2}}\right\rceil-1
 \right)+1
 \right\},
 \qquad M'=dx'.                                      \label{eq:upgrade-combined-floor}
\end{equation}
Across the nine official rows, this requires $45.00\%$--$60.71\%$ more
outputs (\Cref{thm:upgrade-zero-offset-floor}).
\end{theorem}

\begin{proof}
The second and third terms in \eqref{eq:upgrade-combined-floor} are the two
ordered-cap requirements under $m_1'=\lceil x'/d\rceil$.  The term $v$ is
\Cref{prop:upgrade-fresh-square-floor}.  Their maximum is therefore necessary.
Substitution gives the affine-fiber column in the comparison below.
\end{proof}

\paragraph{Zero-offset targets.}
\label{sec:upgrade-zero-offset-repair}

The preceding floor assumes a full-rank affine lift.  A zero-offset attacker
can instead filter chosen-message targets through the quadratic image and use
the full common-column domain.  Removing this square slice requires a stronger
floor.

\begin{proposition}[Zero-offset dimensionally square restriction]
\label{prop:upgrade-zero-offset-square-persistence}
Consider any symmetric redesign for which the zero-offset common-column
relation with $R_0=1$ is defined.  Let $Q_R$ be its quadratic output matrix
and put $\rho_Q=\rank Q_R$.  Public row reduction gives an invertible output
change under which the zero-offset verifier is equivalent to
\begin{equation}
 g_{\rho_Q}(u)=z\in\F_q^{\rho_Q},
 \qquad
 t_{\rm con}=0\in\F_q^{M'-\rho_Q},                     \label{eq:upgrade-general-zero-split}
\end{equation}
where $g_{\rho_Q}$ consists of $\rho_Q$ independent quadratic output
combinations.  For an exactly uniform target, the consistency event in
\eqref{eq:upgrade-general-zero-split} has probability $q^{-(M'-\rho_Q)}$ and,
conditioned on it, $z$ is uniform in $\F_q^{\rho_Q}$.

With zero offsets no reserved line is needed, so use the full public-vinegar
space $V_{\rm vin}\cong A^v$, of base-field dimension $dv$.  If
\begin{equation}
 \rho_Q\le dv,                                          \label{eq:upgrade-zero-square-condition}
\end{equation}
then dimensions alone leave an ordinary $\rho_Q$-dimensional slice
$L\subseteq V_{\rm vin}$ carrying all $\rho_Q$ residual equations: it has
$\rho_Q$ equations in $\rho_Q$ variables.  This dimension equality alone
proves neither root existence nor efficient recovery.  Consequently, ruling
out this chosen-message dimensional configuration requires
\begin{equation}
 \rho_Q>dv.                                             \label{eq:upgrade-zero-square-rank-floor}
\end{equation}
\end{proposition}

\begin{proof}
With zero offsets, the substituted verifier is $Q_Rz_{\rm sym}(u)=t$ and
contains no affine or constant term.  Choose a row basis of $Q_R$ and extend
it to an invertible basis change on the $M'$ output coordinates.  The first
$\rho_Q$ transformed rows are the independent residual quadratics
$g_{\rho_Q}$; the remaining rows vanish identically on the quadratic image and
therefore impose the consistency equations in
\eqref{eq:upgrade-general-zero-split}.  An invertible linear transformation
sends a uniform target to independent uniform residual and consistency
coordinates, proving the distribution statement.  Under
\eqref{eq:upgrade-zero-square-condition}, choose any $\rho_Q$-dimensional
linear subspace of $V_{\rm vin}$ and restrict all $\rho_Q$ residual equations
to an affine coset of it.  This gives the asserted dimensionally square
restriction.
\end{proof}

Symmetry bounds the actual quotient rank by
\begin{equation}
 \rho_Q\le \min\left\{M',\;m_1'\binom{d+1}{2}\right\}.
                                                               \label{eq:upgrade-zero-square-cap}
\end{equation}
Hence making \eqref{eq:upgrade-zero-square-rank-floor} even possible forces
both terms on the right of \eqref{eq:upgrade-zero-square-cap} above $dv$.

\begin{theorem}[Output floor eliminating the zero-offset dimension guarantee]
\label{thm:upgrade-zero-offset-floor}
Retain symmetry, the public vinegar count, and the Version~2.3 coupling
\[
 M'=dx',\qquad m_1'=\left\lceil\frac{x'}d\right\rceil.
\]
Put
\[
 C_d=\binom{d+1}{2},
 \qquad
 m_{\rm sq}=\left\lfloor\frac{dv}{C_d}\right\rfloor+1.
\]
A necessary condition for eliminating the zero-offset chosen-message
dimension guarantee is
\begin{equation}
 x'\ge
 \max\left\{v+1,\;d(m_{\rm sq}-1)+1\right\},
 \qquad M'=dx'.                                        \label{eq:upgrade-zero-offset-floor}
\end{equation}
If the redesign also restores the common-column ordered-label cap
$K_{\rm cc,ord}$, the
right side of \eqref{eq:upgrade-zero-offset-floor} must additionally be
maximized with
\[
  \left\lceil\frac{K_{\rm cc,ord}}d\right\rceil,
 \qquad
  d\left(\left\lceil\frac{K_{\rm cc,ord}}{C_d}\right\rceil-1\right)+1.
\]
For the nine official rows, the zero-offset condition dominates the ordered-cap
terms.  The two repair floors compare as follows.
\begin{center}
\small
\begin{tabular}{@{}c l r r r@{}}
\toprule
&&& \multicolumn{2}{c}{necessary outputs $M'$}\\
\cmidrule(l){4-5}
Level & parameters & current $M$ & affine fiber & zero offset\\
\midrule
I   & $(28,5,4,4)$  & $80$  & $116\;(45.00\%)$ & $180\;(125.00\%)$\\
I   & $(48,16,2,2)$ & $64$  & $96\;(50.00\%)$ & $130\;(103.125\%)$\\
I   & $(28,4,4,5)$  & $80$  & $116\;(45.00\%)$ & $180\;(125.00\%)$\\
III & $(40,7,4,4)$  & $112$ & $180\;(60.71\%)$ & $260\;(132.14\%)$\\
III & $(72,24,2,2)$ & $96$  & $144\;(50.00\%)$ & $194\;(102.08\%)$\\
III & $(38,5,4,5)$  & $100$ & $152\;(52.00\%)$ & $244\;(144.00\%)$\\
V   & $(50,9,4,4)$  & $144$ & $228\;(58.33\%)$ & $324\;(125.00\%)$\\
V   & $(96,32,2,2)$ & $128$ & $192\;(50.00\%)$ & $258\;(101.5625\%)$\\
V   & $(52,6,4,6)$  & $144$ & $228\;(58.33\%)$ & $324\;(125.00\%)$\\
\bottomrule
\end{tabular}
\end{center}
The zero-offset floor requires $101.5625\%$--$144.00\%$ more outputs.
\end{theorem}

\begin{proof}
Condition \eqref{eq:upgrade-zero-square-rank-floor} and the two caps in
\eqref{eq:upgrade-zero-square-cap} require
$dx'>dv$ and $m_1'C_d>dv$.  The first inequality is equivalent to
$x'\ge v+1$.  The second requires $m_1'\ge m_{\rm sq}$, and the least integer
$x'$ with $\lceil x'/d\rceil\ge m_{\rm sq}$ is
$d(m_{\rm sq}-1)+1$.  This proves
\eqref{eq:upgrade-zero-offset-floor}.  The two extra ordered-cap terms are
the same feature and output necessities used in
\Cref{thm:upgrade-combined-floor}.  Direct substitution gives the table.
\end{proof}

The stronger result is a dimension-only theorem.  A redesign below the table
can retain the square system and rely on the consistency factor
$q^{M'-\rho_Q}$, its solving cost, or a changed target interface.  The floor
applies only if the redesign instead removes the square slice while preserving
the common-column architecture and parameter coupling.

\section{Conclusion}\label{sec:conclusion}
The common-column restriction exposes the central weakness of the $q=19$
Version~2.3 construction.  Symmetry collapses its restricted quadratic
coordinates to a symmetric-square quotient.  Conditional on the stated public
ranks, an invertible output change separates this quotient from the affine and
target-consistency constraints.  The result is a quadratic system on a public
affine space.  Every affine restriction preserves the verifier equations, the
rejection rule, and final verification.
This statement concerns the restricted verifier, not the unrestricted
matrix-valued map.

All nine published shapes admit the exact quotient and affine restriction
(\Cref{cor:exact-affine-restriction}).  For $\ell=2$, transverse slices have
full affine coefficient support.  The resulting PXL rows are unit-leading
gate exponents under the ordinary MQ assumptions and the instancewise route
premise (\Cref{thm:pxl-l2-estimates}).  For $\ell=4$, the stronger
$16$-to-$10$
quotient, the fixed-map criterion, and its zero-offset graph realization are
exact for homogeneous channels.  Positive-rank offset compatibility remains
conditional on $\mathsf H_{\rm off}$, whose frequency under the official
distribution is unproved.  Under the instancewise PXL premise, the
support-constrained route crosses Levels III and V but not Level I.  Under the
separate structured PXL premise, the estimates cross all six $\ell=4$ rows
(\Cref{thm:pxl-fixed-map-surjectivity,thm:pxl-all-nine-structured}).  The
artifact contains no production-size solve or forgery
(\Cref{tab:evidence-hierarchy}).

The repair bounds apply only to designs that retain the common-column
architecture and current parameter coupling.  They show that restoring the
ordered-label count does not by itself remove the stronger zero-offset slice
(\Cref{thm:upgrade-combined-floor,thm:upgrade-zero-offset-floor}).

\clearpage
\bibliographystyle{unsrtnat}
\begingroup
\footnotesize
\setlength{\bibsep}{0pt}
\Urlmuskip=0mu plus 1mu\relax
\sloppy
\bibliography{references}
\endgroup

\clearpage
\appendix
\section{Complete PXL model and numerical details}
\label{app:pxl-model}

This appendix states the solver premises and all charged terms behind
\Cref{thm:pxl-l2-estimates}.  PXL first eliminates part of a Macaulay matrix
over a polynomial ring in variables that are fixed later.  It then evaluates
those variables and performs the remaining linear algebra over the base field
~\cite{FurueKudo2024PXL}.

\begin{definition}[Ordinary MQ regularity hypothesis $\mathsf H_{\rm PXL}$]
\label{def:pxl-regularity}
Let $K$ affine quadratic equations in $h$ variables be given.  Choose $k$
variables for polynomial preprocessing and let the remaining $h-k$ variables
be the main variables.  For a degree bound $D$, the hypothesis
$\mathsf H_{\rm PXL}(K,h,k,D)$ consists of the following statements.
\begin{enumerate}[label=(\roman*)]
\item For every assignment to the $k$ fixed variables, the specialized
      system has at most one root over the algebraic closure, counted with
      multiplicity, and its highest-degree parts are semi-regular.
\item The homogeneous quadratic parts in the main variables are
      semi-regular, so their Hilbert series gives the predicted residual
      matrix size.
\item The selected Macaulay rows have the predicted rank, and polynomial
      preprocessing leaves a matrix of the predicted size.
\item Degree $D$ suffices for the second linearization step to produce the
      predicted elimination data.
\end{enumerate}
These assumptions package the affine semi-regularity, projected
semi-regularity, and Macaulay-rank assumptions used in PXL Section~4.2.  They
are not implied by coefficient support and are not proved for the quotient
systems in this paper.
\end{definition}

\begin{definition}[Instancewise route premise $\mathsf H_{\rm route}$]
\label{def:pxl-route}
Fix a public key, retained target, and selected affine slice, and let $f$ be
the resulting restricted system.  In the $\ell=2$ route this is the system
$f_{a,z}$ in \eqref{eq:pxl-restricted-system}.  The premise
$\mathsf H_{\rm route}(f;k,D)$ states that its PXL matrices have
the degrees, dimensions, and ranks specified by
$\mathsf H_{\rm PXL}(K,h,k,D)$.  It also states that, when the system has an
accepted base-field root with full column rank, the manuscript's
specialization and extraction procedure returns at least one such root.
The extraction statement is added by this manuscript and is not part of the
published PXL assumptions.
\end{definition}

For fixed $(K,h,k,D)$, the predicted residual matrix size is
\begin{equation}
 A=\sum_{d=0}^{D}
   \max\left\{
     [t^d](1-t)^{K-h+k}(1+t)^K,
     0
   \right\}.                                         \label{eq:pxl-A}
\end{equation}
The model chooses the least $D\ge2$ such that
\begin{equation}
 [t^D](1-t)^{K-h+k-1}(1+t)^K\le1.                   \label{eq:pxl-degree}
\end{equation}
This is a degree prediction, not a solving-degree theorem for the structured
quotient system.

Put $B=\binom{h-k+D}{D}$.  Section~4.2 of PXL gives five cost components.
We charge all five and aggregate them by addition.  This sum is our
conservative use of those components.  The preprocessing costs are
\begin{align}
 C_1&=B^\omega,                                      \label{eq:pxl-C1}\\
 C_2&=k^2(h-k)B^2,                                   \label{eq:pxl-C2}\\
 C_3&=k^2AB\binom{h+D}{D}.                           \label{eq:pxl-C3}
\end{align}
After preprocessing, evaluation and the second linearization cost
\begin{align}
 C_{\rm fix}&=19^kA^2\binom{k+D}{D},                \label{eq:pxl-Cfix}\\
 C_{\rm lin}&=19^kA^\omega.                         \label{eq:pxl-Clin}
\end{align}
The first three terms are paid once.  The last two charge all $19^k$
assignments to the fixed variables.

Define
\begin{equation}
 C_{\rm PXL}=C_1+C_2+C_3+C_{\rm fix}+C_{\rm lin},
 \qquad C_{\rm target}=2^{32}19^m.                  \label{eq:pxl-kernel-cost}
\end{equation}
Here $m$ is the target-consistency codimension of the route.  The factor
$2^{32}$ is a stipulated unit-leading model convention per retained target.
It is not a measured cost or a circuit-level bound.
Multiplying the field-operation count by the $150$ gates in the released
$\F_{19}$ multiplication circuit gives the multiplication-based gate
normalization
\begin{equation}
 W_{\rm est}=\frac{150C_{\rm PXL}+C_{\rm target}}{p_{\rm dir}}.
                                                               \label{eq:pxl-estimated-work}
\end{equation}
The factor $150$ is not a bound for every operation used by the solver.  If
$C_{\rm aux}$ charges setup,
extraction, control, storage traffic, and final verification for one trial,
the augmented expression is
\begin{equation}
 W_{\rm aug}
 =\frac{150C_{\rm PXL}+C_{\rm target}+C_{\rm aux}}{p_{\rm dir}}.
                                                               \label{eq:pxl-augmented-work}
\end{equation}
Both formulas place target generation inside the retry numerator.  The root
theorem samples a fresh retained target and affine coset in every trial.

The independently checked parameters and matrix sizes are as follows.  The
search is performed separately for each value of $\omega$.
\begin{center}
\scriptsize
\setlength{\tabcolsep}{4pt}
\begin{tabular}{@{}c c r r r r r r@{}}
\toprule
$\omega$ & Level & $k$ & $D$ & $A$ & $\log_2 W_{\rm est}$
& $\log_2 M_{\rm sym}^{\rm proxy}$ & reference exponent\\
\midrule
$2.81$ & I   & 9  & 12 & $99{,}219{,}059$ & 128.833 & 93.337 & 143\\
       & III & 13 & 17 & $2{,}710{,}186{,}080{,}528$ & 187.030 & 137.856 & 207\\
       & V   & 15 & 23 & $547{,}365{,}920{,}911{,}569{,}552$ & 245.042 & 186.447 & 272\\
\midrule
$3$    & I   & 9  & 12 & $99{,}219{,}059$ & 133.667 & 93.337 & 143\\
       & III & 14 & 16 & $962{,}082{,}316{,}942$ & 194.626 & 133.178 & 207\\
       & V   & 19 & 20 & $9{,}230{,}629{,}773{,}094{,}340$ & 255.550 & 172.754 & 272\\
\bottomrule
\end{tabular}
\end{center}
Here
\begin{equation}
 M_{\rm sym}^{\rm proxy}=5B^2\binom{k+D}{D}\quad\text{bits}.  \label{eq:pxl-storage}
\end{equation}
This is a unit-leading dense proxy for one symbolic preprocessing array.  It
is not total state or peak memory.

For Level I, the five charged terms have the following base-two logarithms.
This row makes the aggregation in \eqref{eq:pxl-kernel-cost} reproducible.
\begin{center}
\scriptsize
\begin{tabular}{@{}c r r r r r@{}}
\toprule
$\omega$ & $\log_2 C_1$ & $\log_2 C_2$ & $\log_2 C_3$
& $\log_2 C_{\rm fix}$ & $\log_2 C_{\rm lin}$\\
\midrule
$2.81$ & 102.354 & 84.399 & 109.027 & 109.525 & 112.877\\
$3$    & 109.275 & 84.399 & 109.027 & 109.525 & 117.924\\
\bottomrule
\end{tabular}
\end{center}

\section{Proofs for slice preparation and spectrum bounds}
\label{app:slice-proofs}

This subsection collects the coefficient-level proofs used in
\Cref{sec:sq-preparing,sec:sq-full-domain-spectrum}.  The corresponding
statements and proof ideas remain in the main body.

\begin{proof}[Proof of \Cref{thm:stable-kernel}]
The family in \eqref{eq:Uw} has $m_1d^2$ members.  It is stable under right
multiplication by $A$ because, for every $\zeta\in A$,
\[
 (w^\transpose\Sigma_\xi P_i\Sigma_\eta)\Sigma_\zeta
 =w^\transpose\Sigma_\xi P_i\Sigma_{\eta\zeta},
\]
which is another member of the span.  Every offset cross term contains one
copy of $w$.  If $w$ occurs on the right, symmetry and commutativity move it to
the left.  Thus every row of $F_{R,v}$ lies in $U_w$, and right stability puts
every orbit row there as well.

The complete right orbit in \eqref{eq:orbit-matrix} makes its row space an
$A$-vector space, so $d\mid\rho_F$.  Its first block is $F_{R,v}$, hence its
kernel kills every offset-linear feature.  If $x$ is in the kernel and
$\zeta\in A$, each $S^a\zeta$ is an $\F_q$-linear combination of the chosen
basis powers.  The same orbit rows therefore kill $\Sigma_\zeta x$, proving
that the kernel is $A$-linear.  Rank-nullity over $A$ gives
\eqref{eq:stable-dim}.
\end{proof}

\begin{proof}[Proof of \Cref{lem:reserved-decoupling}]
Every structured offset lies in $W$, so each offset-linear coefficient
contains one copy of $w$.  Symmetry moves that copy to the left when needed.
The coefficient therefore reads only native matrix entries with at least one
index in $W$.  Right multiplication by $A$ preserves this coordinate block,
as do the deterministic orbit and slice constructions.

The native $V'\times V'$ entries come from disjoint XOF symbols and remain
unread.  They are therefore independent of the selected slice in the stated
model.  This argument must stay in the native coordinates because the biased
byte-reduction distribution is not invariant under an arbitrary $A$-linear
basis change.
\end{proof}

\begin{proof}[Proof of \Cref{lem:max-atom-pivot}]
Choose $t$ pivot columns.  Condition on all nonpivot variables.  Expose the
pivot variables in echelon order; at each step at most one value can satisfy
the next independent equation, and that value has probability at most
$\rho_{\rm byte}$.  Multiplication gives the claim.
\end{proof}

\begin{proof}[Proof of \Cref{thm:structural-preflight}]
First consider failure of the affine source rank.  If
$\rank C_{R,v}<M-K$, some nonzero
$\lambda\in\mathcal Q_R^\vee$ annihilates $C_{R,v}$.  In particular,
$(\lambda\circ C_{R,v})|_{V'}=0$.  For fixed $[\lambda]$, this is an affine
system of rank at least $d(v-1)$ in the independent native variables
$\mathbf X_{W,V'}$; \Cref{lem:max-atom-pivot} bounds its probability by
$\rho_{\rm byte}^{d(v-1)}$.  Scalar multiples define the same event, and
$\mathbb P(\mathcal Q_R^\vee)(\F_q)$ has
$(q^{M-K}-1)/(q-1)$ points.  A union bound proves
\eqref{eq:delta-src}.

Next consider failure of the projection to $V'$.  By
\Cref{thm:stable-kernel}, $H^F_{R,v}$ is $A$-linear.  A nonzero intersection
with $W\oplus O_{\rm pub}$ must therefore contain an $A$-line $\mathfrak l$.
For $\mathfrak l=W$, the coefficient-rank hypothesis and
\Cref{lem:max-atom-pivot} give probability at most
$\rho_{\rm byte}^{m_1\binom{d+1}{2}}$.  For every other $A$-line the
analogous bound is $\rho_{\rm byte}^{m_1d^2}$.  The $(1+o)$-dimensional
$A$-space $W\oplus O_{\rm pub}$ contains
$(|A|^{1+o}-1)/(|A|-1)$ lines, exactly one of which is $W$; a union bound
gives \eqref{eq:delta-proj}.  Finally, the kernel of the projection
$H^F_{R,v}\to V'$ is precisely
$H^F_{R,v}\cap(W\oplus O_{\rm pub})$.

A final union bound combines the source-rank and projection failures without
assuming independence.  Direct substitution gives the stated numerical
bound.
\end{proof}

\begin{proof}[Proof of \Cref{lem:sq-full-domain-spectrum}]
Fix $y\ne0$ and $b\ne0$.  Choose a fresh public base form whose corresponding
symmetric-square coefficient $c$ is nonzero.  Writing
$c=\sum_\nu a_\nu\odot b_\nu$, the fresh symmetric-matrix coefficient in the
scalar derivative $b\cdot T_y(x)$ is
\[
 \Xi_{c,\bar y}(x)=
 \sum_\nu\bigl((a_\nu\bar y)\odot(b_\nu x)
 +(b_\nu\bar y)\odot(a_\nu x)\bigr).
\]
We first bound the kernel of this coefficient map.  If $x\notin A\bar y$,
project onto the cross summand between the disjoint lines $A\bar y$ and $Ax$.
After identifying both lines with $A$, this projection is the ordinary
symmetrization of $c$.  It is nonzero because symmetrization is injective on
$\operatorname{Sym}^2(A)$ in odd characteristic.  Hence the kernel is
contained in $A\bar y$, and the coefficient map has rank at least
$\dim_{\F_q}D-d=t$.

The event $b\in\ker T_y^{\mathsf T}$ forces a rank-$t$ affine system in
independent fresh symbols.  Conditioning on the nonpivot symbols and applying
the maximum-atom bound to the pivots proves \eqref{eq:sq-pointwise}.

Finally, the dual-kernel identity gives
\[
 q^{-\rank T_y}
 =q^{-K}|\ker T_y^{\mathsf T}|
 =q^{-K}\left(1+\sum_{b\ne0}
   \mathbf1_{\{b\in\ker T_y^{\mathsf T}\}}\right).
\]
Its expectation is at most
$q^{-K}+(1-q^{-K})\rho_{\rm byte}^t$.  Summing over $y\ne0$ proves
\eqref{eq:sq-eta-expect}.  The excess above the baseline
$(q^h-1)q^{-K}$ is nonnegative, so Markov's inequality proves
\eqref{eq:sq-eta-tail}.  The two stated choices of $D$ have dimensions
$d(v-1)$ and $dv$.
\end{proof}

\paragraph{An unused fixed-skew format check.}
A deterministic alternative is also available:
for the official matrix
\[
 S=\begin{pmatrix}1&2\\2&15\end{pmatrix},
 \qquad
 R_1=9I+10S=\begin{pmatrix}0&1\\1&7\end{pmatrix},
\]
choose offsets with $(v_1)_{b,0}-(v_0)_{b,1}=1$ in every block.  Then each
block has fixed nonzero skew $1$, so every root passes the rejection rule.
This fixed-skew route is a useful separate algebraic check but is not needed
for the reported recovery theorem.


\section{A finite fixed-key spectrum-certificate format}\label{app:fixed-key-spectrum}
The idealized XOF-transcript theorem is not required for a statement about a concrete key.
This appendix gives a finite exact certificate format for the aggregate
spectrum.  No such production-key spectrum certificate is shipped in the
release.

Put $N_U=\dim_{\F_q}U$.  Choose a basis
$y_0,\ldots,y_{h-1}$ of a fixed slice $L$ and write
\[
 T(Y)=\sum_{i=0}^{h-1}Y_iT_{y_i}
 \in\F_q[Y_0,\ldots,Y_{h-1}]^{K\times N_U}.
\]
For $0\le j<h$, use the disjoint normalized projective chart
$Y_0=\cdots=Y_{j-1}=0$, $Y_j=1$, and put
$R_j=\F_q[Y_{j+1},\ldots,Y_{h-1}]$.  Let $T^{(j)}$ be the resulting matrix.
For $t\ge0$, define
\begin{equation}
 I_{j,t}=\left\langle
   Y_i^q-Y_i\ (i>j),\quad
   \text{all }(t+1)\times(t+1)\text{ minors of }T^{(j)}
 \right\rangle\subseteq R_j.                        \label{eq:histogram-ideal}
\end{equation}
Put $n_{j,\le t}=\dim_{\F_q}(R_j/I_{j,t})$,
$n_{j,\le-1}=0$, and
\begin{equation}
  N_t=\sum_{j=0}^{h-1}(n_{j,\le t}-n_{j,\le t-1}). \label{eq:histogram-count}
\end{equation}

\begin{proposition}[Projective rank-histogram certificate]
\label{prop:projective-histogram}
The integer $N_t$ is exactly the number of projective directions
$[y]\in\mathbb P(L)(\F_q)$ for which $\rank T_y=t$.  Consequently
\begin{equation}
  \bar\eta_L^U=(q-1)\sum_tN_tq^{-t}.                 \label{eq:histogram-eta}
\end{equation}
Reduced Gr\"obner bases with their standard monomial sets, or polynomial
reduction transcripts proving the unit-ideal and quotient-dimension claims,
form machine-checkable fixed-key certificates.
\end{proposition}

\begin{proof}
The normalized charts contain exactly one representative of every projective
$\F_q$-direction.  The field equations identify
\[
 R_j/\langle Y_i^q-Y_i:i>j\rangle
 \cong\prod_{a\in\F_q^{h-j-1}}\F_q
\]
by evaluation.  Every ideal in this product of fields is radical, and its
quotient dimension is the number of coordinates on which all generators
vanish.  The minors in \eqref{eq:histogram-ideal} vanish exactly when
$\rank T^{(j)}(a)\le t$.  Thus $n_{j,\le t}$ counts rank-at-most-$t$ points
in chart $j$, and differencing gives \eqref{eq:histogram-count}.  Every
projective direction has $q-1$ nonzero scalar representatives and rank is
invariant under base-field scaling, proving \eqref{eq:histogram-eta}.
\end{proof}

A stronger certificate that $\rank T_y\ge t_0+1$ for every nonzero $y$
consists of unit-ideal transcripts for all $I_{j,t_0}$.  The complete histogram is more
flexible: a small exceptional low-rank locus can be charged exactly instead of
controlling the whole pencil by its minimum rank.

\section{Exact \texorpdfstring{$\ell=4$}{ell=4} PXL ledger rows}
\label{app:pxl-ell4-ledger}
The following tables record the exact inputs and outputs in the final PXL
ledger.  The generator and checker verify the formulas independently.  The
support-constrained rows condition on
$\mathsf H_{\rm off}$ and $\mathsf H_{\rm route}$.  The structured rows
condition on $\mathsf H_{\rm struct}$ and do not assert ordinary coefficient
support.  Together with the $\ell=2$ table in \Cref{app:pxl-model}, these are
the 30 rows checked by the artifact.

\begin{table}[H]
\centering
\caption{Exact support-constrained $\ell=4$ PXL ledger rows.}
\label{tab:pxl-l4-supported-ledger}
\scriptsize
\setlength{\tabcolsep}{3pt}
\begin{tabular}{@{}c l r r r r r r@{}}
\toprule
$\omega$ & parameters & $h$ & $k$ & $D$ & $A$ & $\log_2 W_{\rm est}$ & $\log_2 M_{\rm sym}^{\rm proxy}$\\
\midrule
$2.81$ & $(28,5,4,4)$ & 35 & 7 & 7 & $60{,}179$ & 145.867751 & 59.429\\
       & $(28,4,4,5)$ & 35 & 7 & 7 & $60{,}179$ & 145.867386 & 59.429\\
       & $(40,7,4,4)$ & 48 & 8 & 9 & $4{,}927{,}544$ & 197.338396 & 78.763\\
       & $(38,5,4,5)$ & 46 & 8 & 8 & $1{,}606{,}800$ & 201.265361 & 71.892\\
       & $(50,9,4,4)$ & 59 & 9 & 10 & $108{,}000{,}036$ & 252.202579 & 91.085\\
       & $(52,6,4,6)$ & 61 & 9 & 10 & $461{,}010{,}165$ & 249.538960 & 92.109\\
\midrule
$3$    & $(28,5,4,4)$ & 35 & 7 & 7 & $60{,}179$ & 148.400390 & 59.429\\
       & $(28,4,4,5)$ & 35 & 7 & 7 & $60{,}179$ & 148.400025 & 59.429\\
       & $(40,7,4,4)$ & 49 & 9 & 9 & $4{,}927{,}544$ & 201.379285 & 79.763\\
       & $(38,5,4,5)$ & 46 & 8 & 8 & $1{,}606{,}800$ & 205.022213 & 71.892\\
       & $(50,9,4,4)$ & 60 & 10 & 10 & $108{,}000{,}036$ & 257.208548 & 92.085\\
       & $(52,6,4,6)$ & 61 & 9 & 10 & $461{,}010{,}165$ & 254.991164 & 92.109\\
\bottomrule
\end{tabular}
\end{table}

\begin{table}[H]
\centering
\caption{Exact structured $\ell=4$ PXL ledger rows.}
\label{tab:pxl-l4-structured-ledger}
\scriptsize
\setlength{\tabcolsep}{3pt}
\begin{tabular}{@{}c l r r r r r r@{}}
\toprule
$\omega$ & parameters & $h$ & $k$ & $D$ & $A$ & $\log_2 W_{\rm est}$ & $\log_2 M_{\rm sym}^{\rm proxy}$\\
\midrule
$2.81$ & $(28,5,4,4)$ & 48 & 10 & 12 & $125{,}101{,}086$ & 134.037961 & 95.266\\
       & $(28,4,4,5)$ & 48 & 10 & 12 & $125{,}101{,}086$ & 134.037595 & 95.266\\
       & $(40,7,4,4)$ & 68 & 12 & 17 & $2{,}232{,}564{,}665{,}436$ & 182.006674 & 135.896\\
       & $(38,5,4,5)$ & 68 & 12 & 17 & $2{,}232{,}564{,}665{,}436$ & 182.006152 & 135.896\\
       & $(50,9,4,4)$ & 88 & 15 & 21 & $14{,}990{,}431{,}104{,}743{,}028$ & 230.444951 & 172.076\\
       & $(52,6,4,6)$ & 88 & 15 & 21 & $14{,}990{,}431{,}104{,}743{,}028$ & 230.444296 & 172.076\\
\midrule
$3$    & $(28,5,4,4)$ & 48 & 10 & 12 & $125{,}101{,}086$ & 138.911992 & 95.266\\
       & $(28,4,4,5)$ & 48 & 10 & 12 & $125{,}101{,}086$ & 138.911626 & 95.266\\
       & $(40,7,4,4)$ & 68 & 13 & 16 & $811{,}505{,}434{,}656$ & 189.638021 & 131.353\\
       & $(38,5,4,5)$ & 68 & 13 & 16 & $811{,}505{,}434{,}656$ & 189.637499 & 131.353\\
       & $(50,9,4,4)$ & 88 & 18 & 19 & $706{,}353{,}332{,}229{,}810$ & 240.174424 & 162.956\\
       & $(52,6,4,6)$ & 88 & 18 & 19 & $706{,}353{,}332{,}229{,}810$ & 240.173769 & 162.956\\
\bottomrule
\end{tabular}
\end{table}

\section{Arithmetic-circuit certificate details}\label{app:circuits}
The scalar multiplier consumes two little-endian five-bit canonical integers
in $\{0,\ldots,18\}$ and returns the canonical five-bit representation of
their product modulo $19$.  Its $150$ gates are two-input AND, XOR, and XNOR
gates with free constants, wires, fan-out, and structural sharing.  The
released netlist is tested on all $361$ valid input pairs.  The SHA-256 digest
of the canonical compact JSON payload containing exactly
\texttt{ninputs}, \texttt{gates}, and \texttt{outputs} is
\begin{center}
\ttfamily\small
 e1708a32c1475ad14295c2b34fed976fb09577b127810d29bf237db18aabc131
\end{center}
for the emitted balanced-netlist representation.  This is an internal payload
identifier, not the bytewise digest of the surrounding ledger file.

The tower recurrences are
\[
 \F_{19^2}=\F_{19}[u]/(u^2+1),
 \qquad
 \F_{19^4}=\F_{19^2}[v]/(v^2-(1+u)).
\]
Three-product Karatsuba at each level gives separately checked recurrence
counts $692$ and $2628$.  The artifact verifies the tower identities, the
basis determinant, and the recurrence arithmetic; it does not emit or
exhaustively test complete extension-field netlists.  These extension-field
counts are not used in the PXL tables.

The scalar multiplication circuit supplies the multiplication-based
normalization used in the PXL tables.  The circuits are not lower bounds and
do not price the whole solver.  Additions,
inversions, factorization, bulk evaluation/interpolation transforms,
filtering, control, indexing, communication, and memory are outside the
reported $W_{\rm est}$ and belong to $C_{\rm aux}$ if charged.

\clearpage
\section{Reproduction and package contents}\label{app:artifact}
The source package is self-contained.  Its root contains the complete paper
source, bibliography, compiled PDF, build instructions, and hashes.  The
\texttt{artifact/} directory contains:
\begin{itemize}
\item the final PXL generator and pinned JSON ledger for all 30 rows,
      comprising six $\ell=2$, twelve support-constrained $\ell=4$, and
      twelve structured $\ell=4$ rows;
\item an aggregate verifier that reruns the PXL checker and the retained
      formula, target-sampling, dimension, and repair-floor checks;
\item the complete $150$-gate scalar multiplier netlist and tower checks;
\item a public-key correspondence harness for one official Level-I KAT; and
\item a non-planted zero-offset end-to-end composition test at the unofficial
      shape $(2,1,2,2)$, with fresh-key forgery and rank stress checks.
\end{itemize}

A clean reproduction is:
\begin{verbatim}
sha256sum -c SHA256SUMS
make regenerate
make verify
make
\end{verbatim}
The PXL ledger records $\mathsf H_{\rm off}$ only for the
support-constrained $\ell=4$ route and $\mathsf H_{\rm struct}$ only for the
structured route.  It does not assign either premise an empirical
probability.  The checkers are deterministic, require Python~3.10 or newer
plus NumPy for the two correspondence tests, and do not require network
access.  The paper build uses only the files in the package.  The supplied
\texttt{SHA256SUMS} records the release-critical files listed in the manifest.
Generated TeX build files are excluded.

\end{document}
